\section{Sparse equations for singular moduli}\label{sec:smd}
% A ragged paragraph listing delivered certificate files (long paths do not hyphenate).
\providecommand{\smdcert}[1]{\par\smallskip{\small\raggedright\noindent
  \emph{Certificates:} #1\par}\smallskip}

A \emph{singular modulus} is a value $\alpha=j(\tau)$ of the modular invariant at an
imaginary quadratic point $\tau$ of the upper half-plane. The corresponding elliptic curve
has complex multiplication by an imaginary quadratic order; we write $\Delta=-D<0$ for the
discriminant of that order, $\Delta=f^2d_K$ with $d_K$ the fundamental discriminant of
$K=\Q(\sqrt\Delta)$ and $f$ the conductor, $h=h(\Delta)$ for the class number of the order,
and $H_\Delta\in\Z[X]$ for the monic ring class polynomial. The conjugates of $\alpha$ over
$\Q$ are exactly the $h$ singular moduli of discriminant $\Delta$, one for each primitive
reduced positive definite form $[a,b,c]$ of discriminant $\Delta$, namely
$j(\tau)$ with $\tau=(-b+i\sqrt D)/(2a)$; they are the roots of $H_\Delta$ and are
algebraic integers.

The \emph{support} of a polynomial is its set of nonzero coefficients. Define
\[
 s(\Delta)=\min\set{\#\supp P : 0\ne P\in\Q[X],\ P(\alpha)=0}.
\]
All conjugates have the same minimum, so $s$ depends only on $\Delta$. Equivalently,
$s(\Delta)$ is the least number of terms of a nonzero multiple of $H_\Delta$. The degree of
$P$ and the size of its coefficients are unrestricted, and the class polynomial itself gives
\begin{equation}\label{smd:eq:trivial}
 s(\Delta)\le h(\Delta)+1 .
\end{equation}
The statement $s(\Delta)>t$ says that \emph{any} $t$ distinct nonnegative powers of
$\alpha$ are linearly independent over $\Q$.

Riffaut's conjecture asserts that a singular modulus of degree at least three is not a root
of a trinomial, that is, $s(\Delta)\ge4$ whenever $h(\Delta)\ge3$. Bilu, Luca and
Pizarro-Madariaga proved it assuming GRH for Dirichlet $L$-functions
\cite{Bilu2022Trinomials}, a hypothesis which \cref{src:thm:dgrh} supplies. This section
determines $s(\Delta)$ in every degree up to five and classifies the equations with at most
five terms without any hypothesis on zeros (\cref{smd:thm:five}); it shows that
$s(\Delta)$ has the order of magnitude of the class number up to a power of
$\log D$ when Dirichlet GRH is used (\cref{smd:thm:growth}), with weaker unconditional
bounds (\cref{smd:thm:uncond,smd:thm:nearfull}); and it exhibits infinite families of orders
for which the trivial bound \eqref{smd:eq:trivial} is an equality
(\cref{smd:thm:dyadic,smd:thm:gaussian}), among them fourteen orders of degrees six and seven
(\cref{smd:thm:seven}). It also determines the exact number of roots that a polynomial with
at most six terms can have in a deep multiplicative $p$-adic disk
(\cref{smd:thm:rootcount,smd:thm:sixroot}).

\subsection{Coefficients in the CM field}
{\raggedright\sources{grh_cm_prime_applications; grh_cm_four_terms_and_prime_bias; grh_sparse_singular_moduli}}

\begin{proposition}[Coefficients in the CM field]\label[proposition]{smd:prop:cmfield}
Let $\alpha$ be a singular modulus of discriminant $\Delta$ and $K=\Q(\sqrt\Delta)$. The
minimum support size of a nonzero polynomial vanishing at $\alpha$ is exactly $s(\Delta)$
when the coefficients range over $K$. The same holds for any coefficient field linearly
disjoint from $\Q(\alpha)$ over $\Q$.
\end{proposition}

\begin{proof}
The ring class field theorem gives $[K(\alpha):K]=h(\Delta)=[\Q(\alpha):\Q]$
\cite{Cox2013Primes}. Consequently the natural map $K\otimes_\Q\Q(\alpha)\to K(\alpha)$ is
an isomorphism, and the rational basis $1,\alpha,\dots,\alpha^{h-1}$ of $\Q(\alpha)$ remains
a basis of $K(\alpha)$ over $K$. Every finite set of powers of $\alpha$ therefore has the
same rank over the two fields. Conversely a rational relation is also a relation over $K$.
The argument uses only linear disjointness.
\end{proof}

All results below about $s(\Delta)$ therefore hold verbatim for polynomials with
coefficients in the corresponding imaginary quadratic field.

\subsection{Sparse polynomials in a multiplicative \texorpdfstring{$p$}{p}-adic disk}\label{smd:sec:disk}
{\raggedright\sources{grh_cm_sparse_equations; grh_cm_four_terms_and_prime_bias; grh_sparse_singular_moduli}}

Normalize $v_p$ by $v_p(p)=1$ on $\overline{\Q}_p$. In each argument a finite extension of
$\Q_p$ contains all points and scaling constants used; $\OO$ denotes its valuation ring. A
\emph{multiplicative disk of relative depth $r$} about $a\ne0$ is the set of $x$ with
$v_p(x/a-1)\ge r$. For rational $r>1/(p-1)$ put
\[
 w_{p,r}(n)=rn-v_p(n!)\qquad(n\ge0),
\]
so that $w_{p,r}(0)=0$ and $w_{p,r}(n)\to\infty$ by Legendre's formula
$v_p(n!)=(n-S_p(n))/(p-1)$, where $S_p(n)$ is the sum of the base-$p$ digits of $n$.

\begin{lemma}[Local root count]\label[lemma]{smd:lem:disk}
Let $P(X)=\sum_{j=1}^t c_jX^{n_j}$ have distinct nonnegative integer exponents and
nonzero coefficients in a finite extension of $\Q_p$, let $a\ne0$, and let $r>1/(p-1)$ be
rational. The number of distinct roots of $P$ in the multiplicative disk of relative depth
$r$ about $a$ is at most
\[
 B_p(t,r)=\max\set{k\ge0 : w_{p,r}(k)\le\max_{0\le i<t}w_{p,r}(i)} .
\]
\end{lemma}

\begin{proof}
Enlarge the field to contain the roots, the centre and an element $b$ with $v_p(b)=r$.
The substitution $X=a\exp(bY)$, with inverse $Y=b^{-1}\log(X/a)$, identifies the disk
with $\set{v_p(Y)\ge0}$. For the moments $M_k=\sum_jc_ja^{n_j}n_j^k$ (with $n_j^0=1$),
\[
 P(a\exp(bY))=\sum_{k\ge0}M_k\,\frac{b^k}{k!}\,Y^k .
\]
The first $t$ moments are not all zero, by the ordinary Vandermonde determinant in the
distinct integers $n_j$. All later moments obey the monic integral recurrence with
characteristic polynomial $\prod_j(T-n_j)\in\Z[T]$, so
$v_p(M_k)\ge m:=\min_{i<t}v_p(M_i)$ for every $k$ by the ultrametric inequality. Some
$i<t$ attains $m$, and the coefficient of $Y^i$ then has valuation $m+w_{p,r}(i)$. When
$w_{p,r}(k)>\max_{i<t}w_{p,r}(i)$, the coefficient of $Y^k$ has strictly larger valuation.
The series converges on the closed unit disk because $w_{p,r}(k)\to\infty$. Strassmann's
theorem bounds the number of zeros by the last index of a coefficient of largest absolute
value, which is at most $B_p(t,r)$.
\end{proof}

The lemma allows rational depths in ramified fields, algebraic centres, and disks centred
at a nonunit after a common scaling. It is related to Lenstra's local fewnomial bounds and
to Proposition 3.8 of Avenda\~no and Krick, which gives the bound $t-1$ in a principal-unit
disk when $p>e+t-1$ for ramification index $e$
\cite{Lenstra1999Factorization,Avendano2011Sharp}; retaining the individual factorial
valuations handles small residue characteristic and values of $t$ much larger than $p$.
Since $B_p(t,r)$ is nondecreasing in $t$, a bound for $t$ terms also covers polynomials
with fewer terms. Two consequences are used repeatedly.

\begin{corollary}[Depth thresholds]\label[corollary]{smd:cor:threshold}
Let $t\ge2$ and let $\kappa\ge1/(p-1)$ satisfy
\begin{equation}\label{smd:eq:factorial}
 v_p\bigl((i+d)!/i!\bigr)\le d\kappa\qquad(0\le i<t,\ d\ge1).
\end{equation}
If $t$ distinct numbers $z_1,\dots,z_t$ satisfy $v_p(z_l/a-1)>\kappa$ for one $a\ne0$,
then $\det(z_l^{n_j})_{l,j\le t}\ne0$ for every choice of distinct nonnegative integers
$n_1,\dots,n_t$; equivalently, no nonzero polynomial with at most $t$ terms vanishes at all
the $z_l$. Condition \eqref{smd:eq:factorial} holds for
\[
 \kappa(p,t)=\frac1{p-t+1}\quad(p\ge t),\qquad
 \kappa(2,t)=2,\quad \kappa(3,t)=1\quad(t\le5).
\]
\end{corollary}

\begin{proof}
Choose a rational $r$ with $\kappa<r\le\min_lv_p(z_l/a-1)$, so that $r>1/(p-1)$. For $i<t\le n$, condition \eqref{smd:eq:factorial} with $d=n-i$ gives
$w_{p,r}(n)-w_{p,r}(i)=r(n-i)-v_p(n!/i!)>0$. Hence $B_p(t,r)<t$, and
\cref{smd:lem:disk} shows that a polynomial with at most $t$ terms has fewer than $t$ roots
in the disk. A vanishing determinant would give such a polynomial with exponents among the
$n_j$.

For $p\ge t$ the left side of \eqref{smd:eq:factorial} is zero when $i+d<p$. Otherwise,
Legendre's formula gives $v_p((i+d)!)\le(i+d-1)/(p-1)$, and
$(i+d-1)/(p-1)\le d/(p-i)$ for $i\ge1$ and $d\ge p-i$, because the difference of the
cross-products is $(i-1)(p-i-d)\le0$. For $i=0$ use $v_p(d!)\le d/(p-1)$. Since
$i\le t-1$, each bound is at most $d/(p-t+1)$. For $p\in\set{2,3}$, $i\le4$ and $d\ge3$,
the same formula gives $(i+d-1)/(p-1)\le2d$ at $p=2$ and $\le d$ at $p=3$; for $d=1,2$
one checks the products of the factors $i+1,i+2\le6$ directly.
\end{proof}

\begin{corollary}[Depth one in an unramified disk]\label[corollary]{smd:cor:coarse}
Let $p$ be odd and let $P$ have $t$ nonzero terms. A multiplicative disk of relative depth
$1$ contains at most
\begin{equation}\label{smd:eq:coarse}
 \floor*{\frac{p-1}{p-2}\,(t-1)}
\end{equation}
distinct roots of $P$. More precisely, if it contains $N$ distinct roots, then
\begin{equation}\label{smd:eq:digits}
 N\le t-1+\frac{p-1}{p-2}\bigl(1+\log_pt\bigr).
\end{equation}
\end{corollary}

\begin{proof}
Here $w_{p,1}(n)=n-v_p(n!)=\bigl((p-2)n+S_p(n)\bigr)/(p-1)$. Since
$w_{p,1}(n)\ge n(p-2)/(p-1)$ and $\max_{i<t}w_{p,1}(i)\le t-1$, the first bound follows
from \cref{smd:lem:disk}. For the second, the lemma gives some $k\ge N$ and some $i<t$ with
$(p-2)k+S_p(k)\le(p-2)i+S_p(i)$. Use $S_p(k)\ge1$ and
$S_p(i)\le(p-1)(1+\log_pt)$.
\end{proof}

The bound \eqref{smd:eq:coarse} holds even when $t$ greatly exceeds $p$.

\begin{lemma}[Counting conjugates in a disk]\label[lemma]{smd:lem:newton}
Let $H$ be a monic separable polynomial with $p$-adically integral coefficients, let
$v\ge0$ be an integer and $u$ a rational integer prime to $p$, and expand
\[
 H\bigl(p^v(u+Z)\bigr)=\sum_iB_iZ^i .
\]
For every rational $\lambda$, the least index minimizing $v_p(B_i)+i\lambda$ (omitting zero
coefficients) equals the number of roots $\alpha$ of $H$ with
$v_p\bigl(\alpha/(p^vu)-1\bigr)>\lambda$.
\end{lemma}

\begin{proof}
The roots $z$ of the displayed polynomial in $Z$ are the numbers $\alpha/p^v-u$, and
$v_p(\alpha/(p^vu)-1)=v_p(z)$. Choose $\pi$ with $v_p(\pi)=\lambda$, substitute
$Z=\pi Y$, and divide by a coefficient of smallest valuation. Factor into linear factors
in a splitting field and reduce. A root with $v_p(z)>\lambda$ contributes a factor $Y$,
a root with $v_p(z)=\lambda$ a factor $Y$ minus a nonzero residue, and a root with
$v_p(z)<\lambda$ a nonzero constant after normalization. The order of vanishing at $Y=0$
of the reduction therefore counts exactly the roots strictly inside, and it equals the
least minimizing index. This is the Newton polygon rule, including ties.
\end{proof}

Rational noninteger radii thus require no computation in a ramified field: the test uses
only the integers $v_p(B_i)$.

\begin{example}[Two certificates by hand]\label[example]{smd:ex:hand}
For $\Delta=-55$ the valuations at $3$ of the five coefficients of $H_{-55}(27(1+Z))$,
constant term first, are $(20,18,17,14,12)$. With $\lambda=1$ the weights
$v_3(B_i)+i$ are $(20,19,19,17,16)$, with unique minimum at index $4$. All four roots
have relative depth greater than $1=\kappa(3,4)$, so $s(-55)=5$ by
\cref{smd:cor:threshold} and \eqref{smd:eq:trivial}.

For $\Delta=-68$ expand $2^{-24}H_{-68}(2^6(3+Y))$. From constant to leading term the
$2$-adic valuations are $13,10,6,5,0$. With $\lambda=5/2$ the unique minimizing index of
$v_2(B_i)+5i/2$ is $4$, so all four conjugates lie in the disk of relative depth $5/2$
about $2^6\cdot3$. For $i\le3$ the maximum of $w_{2,5/2}(i)=5i/2-v_2(i!)$ is $13/2$; at
$n=4$ the weight is $7$, and for $n\ge5$ it is at least $3n/2+1$. Hence
$B_2(4,5/2)=3$ and $s(-68)=5$.
\end{example}

\subsection{Exact root counts for at most six terms}\label{smd:sec:exact}
{\raggedright\sources{padic-fewnomial-root-count; padic-six-term}}

When $t\le p$ the factorial bound $B_p(t,r)$ of \cref{smd:lem:disk} can be far from the
truth at shallow depths. For a prime $p$, an integer $t\ge2$ and a real $r>1/(p-1)$, let
$N_p(t,r)$ be the largest number of distinct roots that a nonzero polynomial with at most
$t$ terms can have in a closed multiplicative disk $v_p(x/a-1)\ge r$ about a nonzero centre
$a$, the roots being allowed in any finite extension of $\Qp$.

\begin{theorem}[Two to five terms; \stU]\label{smd:thm:rootcount}
Let $2\le t\le\min\set{5,p}$ and $r>1/(p-1)$. Then $N_p(2,r)=1$,
\[
 N_p(t,r)=\begin{cases}p,& r\le1/(p-t+1),\\ t-1,& r>1/(p-t+1),\end{cases}\qquad(t=3,4),
\]
and
\[
 N_p(5,r)=\begin{cases}p+1,& r\le1/(p-3),\\ p,& 1/(p-3)<r\le1/(p-4),\\ 4,& r>1/(p-4).\end{cases}
\]
The upper bounds hold for roots counted with multiplicity, about every centre. In every
range the maximum is attained, about the centre $1$, by a polynomial with integer
coefficients that is irreducible over $\Q$. Moreover, for all $2\le t\le p$ and all
$r>1/(p-1)$,
\begin{equation}\label{smd:eq:uniform}
 N_p(t,r)\le\max\set{t-1,\ p(t-2)} .
\end{equation}
\end{theorem}

\begin{theorem}[Six terms; \stU]\label{smd:thm:sixroot}
Let $p\ge7$ and $r>1/(p-1)$. Then
\[
 N_p(6,r)=\begin{cases}
 2p,& r\le1/(p-2),\\
 p+1,& 1/(p-2)<r\le1/(p-4),\\
 p,& 1/(p-4)<r\le1/(p-5),\\
 5,& r>1/(p-5),
 \end{cases}
\]
with the same statements on multiplicity and on rational irreducible witnesses about $1$
as in \cref{smd:thm:rootcount}. At $p=7$ the four values are $14,8,7,5$.
\end{theorem}

\begin{example}[Comparison with the factorial bound]\label[example]{smd:ex:factorial}
Let $p=t=5$ and $r=1/4+10^{-9}$. Then $\max_{i<5}w_{5,r}(i)=w_{5,r}(4)=1+4\cdot10^{-9}$,
while $k=2\cdot5^{12}=488{,}281{,}250$ has
$w_{5,r}(k)=(r-1/4)k+S_5(k)/4=0.48828125+0.5<w_{5,r}(4)$. Hence
$B_5(5,r)\ge488{,}281{,}250$, whereas $N_5(5,r)=6$ by \cref{smd:thm:rootcount}.
\end{example}

The general estimate $(n-2)p$ of order $n$ over $\Qp$, for $p>n\ge3$, is classical
\cite{vanderPoortenRumely1987Zeros}, and \cite{Avendano2011Sharp} bounds the roots
in a fixed local field; \eqref{smd:eq:uniform} has the same shape, while the two theorems give
the exact dependence on the depth, with roots in arbitrary finite extensions. The proofs
refine the moment argument of \cref{smd:lem:disk} by reducing the moment recurrence modulo
$p$.

\paragraph{The residue algebra.}
Let $2\le t\le p$, let $P=\sum_{j=1}^tc_jX^{n_j}$ be nonzero with $t$ distinct nonnegative
exponents (zero coefficients are allowed, which covers polynomials with fewer terms), and
let $a\ne0$ and $r>1/(p-1)$ be rational. As in the proof of \cref{smd:lem:disk}, the
coefficient of $Y^k$ in $P(a\exp(bY))$, $v_p(b)=r$, is $M_kb^k/k!$ with
$M_k=\sum_jc_ja^{n_j}n_j^k$, and its valuation is $v_p(M_k)+w(k)$, where
\[
 w(k)=w_{p,r}(k)=\Bigl(r-\frac1{p-1}\Bigr)k+\frac{S_p(k)}{p-1}.
\]
By the Weierstrass preparation theorem the number of roots in the disk, counted with
multiplicity, is the largest index $k$ at which the coefficient valuation attains its
minimum $\mu$; the substitution is an analytic bijection with nonvanishing derivative. Call
an index \emph{excluded} if it does not attain $\mu$. Divide all moments by one of least
valuation among $M_0,\dots,M_{t-1}$, and write $m_k$ for the quotients; by the integral
recurrence all $m_k$ are integral and some $m_{i_0}$ with $i_0<t$ is a unit. Thus
$m_k=L(T^k)$ for the functional $L(f)=\sum_jc_ja^{n_j}f(n_j)/M_{i_0}$ on
$\Lambda=\Z[T]/\prod_j(T-n_j)$, which is integral. Put
\[
 \mathcal R=\Lambda/p\Lambda=\F_p[T]\Big/\prod_j(T-\bar n_j),\qquad U=T^p,\qquad V=T-U .
\]
On the factor of $\mathcal R$ at a residue $c$, of multiplicity at most $t\le p$, one has
$U=c^p+(T-c)^p=c$ and $V=T-c$; hence $T^{p^j}=U$ for all $j\ge1$. The \emph{profile} is the
multiset of multiplicities of the residues $\bar n_j$.

\begin{lemma}[Excluded indices]\label[lemma]{smd:lem:excluded}
In this situation:
\begin{enumerate}[label=\textup{(\alph*)}]
\item an index $k\ge t$ with $w(k)>w(t-1)$ is excluded; in particular every $k\ge t$ with
$S_p(k)\ge t-1$ is excluded, and if $r>1/(p-t+1)$ every $k\ge t$ is excluded;
\item an index $k\ge t$ is excluded if the class of $T^k$ in $\mathcal R$ lies in the
$\F_p$-span of classes $T^\ell$ with $w(\ell)<w(k)$, for instance with $\ell<k$ and
$S_p(\ell)\le S_p(k)$;
\item if $k=\sum_ja_jp^j\ge t$ in base $p$ is not excluded, then $k=a_0+pa_1$ with
$a_0+a_1\le t-2$.
\end{enumerate}
Consequently $N_p(t,r)\le\max\set{t-1,p(t-2)}$.
\end{lemma}

\begin{proof}
For $i<t\le p$ one has $v_p(i!)=0$, so $w(i)=ri$ increases on $\set{0,\dots,t-1}$ and
$\mu\le w(i_0)\le w(t-1)$, after the normalization. This proves the first claim of (a). If
$k\ge t$ and $S_p(k)\ge t-1$, then $w(k)-w(t-1)\ge(r-1/(p-1))(k-t+1)>0$. If
$r>1/(p-t+1)$ and $t\le k<p$, then $w(k)=rk>w(t-1)$; for $k\ge p$, Legendre's formula gives
$v_p(k!)\le(k-1)/(p-1)\le(k-t+1)/(p-t+1)<r(k-t+1)$, the middle inequality being equivalent
to $(k-p)(t-2)\ge0$, and again $w(k)>w(t-1)$.

(b) For $i<t\le k$, $w(k)-w(i)>(S_p(k)-i)/(p-1)\ge(1-(p-1))/(p-1)>-1$, so $w(i)-w(k)<1$.
If $v_p(m_k)\ge1$, the coefficient of index $k$ has valuation at least
$1+w(k)>w(i_0)\ge\mu$ in normalized units, and $k$ is excluded. If $v_p(m_k)<1$, lift the
relation $T^k\equiv\sum_\ell\lambda_\ell T^\ell$ to integers; since $\Lambda$ is a free
$\Z$-module the difference lies in $p\Lambda$, so $m_k-\sum_\ell\lambda_\ell m_\ell\in p\OO$
by integrality of $L$. Some $\ell$ then has $v_p(m_\ell)\le v_p(m_k)$, and its coefficient
has strictly smaller valuation because $w(\ell)<w(k)$. The parenthetical case follows from
$w(k)-w(\ell)=(r-1/(p-1))(k-\ell)+(S_p(k)-S_p(\ell))/(p-1)$. This argument is valid in
ramified extensions, because reduction is modulo the ideal $p\OO$.

(c) By (a), $S_p(k)\le t-2\le p-2$. Put $k'=a_0+p\sum_{j\ge1}a_j$. Then $T^k$ and $T^{k'}$
have the same class $T^{a_0}U^{\sum_{j\ge1}a_j}$ in $\mathcal R$, $k'\le k$, and
$S_p(k')=S_p(k)$ since $\sum_{j\ge1}a_j<p$. If $k'<k$, part (b) excludes $k$.

Finally the largest non-excluded index is either at most $t-1$ or of the form $a_0+pa_1$
with $a_0+a_1\le t-2$, hence at most $(t-2-a_1)+pa_1\le p(t-2)$. For an irrational depth
apply the bound at a slightly smaller rational depth; the disk only grows.
\end{proof}

Write $T^iU^j$ for the class of $T^{i+pj}$. By \cref{smd:lem:excluded}(c) the indices $k\ge t$
that remain are those of the classes $T^iU^j$ with $j\ge1$ and $i+j\le t-2$, and the index of
$T^iU^j$ is $i+pj$. In the span arguments below every class is expressed through classes of
smaller index and of digit sum at most its own, so \cref{smd:lem:excluded}(b) applies. We use
three identities in $\mathcal R$: $(U-c)(U-c')=0$ if there are at most two residues $c,c'$;
$V^2=0$ if all multiplicities are at most two; and $(U-c)V=0$ if $c$ is the only repeated
residue. Note that $(U-c)V=0$ gives $U^2=TU-cT+cU$.

\begin{proof}[Proof of the upper bounds in \cref{smd:thm:rootcount}]
By \cref{smd:lem:excluded}(a) all bounds equal $t-1$ when $r>1/(p-t+1)$; this gives $t=2$
and the last range in each case. For $t=3$ the only remaining class is $U$, of index $p$.

For $t=4$ the remaining classes are $U$, $TU$ and $U^2$; we show that $TU$ and $U^2$ lie in
the span of $1,T,T^2,U$. If all multiplicities are at most two, $V^2=0$ gives
$2TU=T^2+U^2$, and $U^2$ is affine in $U$ when there are at most two residues; with at least
three residues there is at most one repeat $c$, and combining $(U-c)V=0$ with $V^2=0$ gives
$U^2=T^2+2cU-2cT$ (if there is no repeat, $U=T$). In profile $(3,1)$, with $c$ repeated and
$c'$ simple, $TU=cT+c'(U-c)$, as one checks on both factors, and $U^2$ is affine in $U$. In
profile $(4)$ the class $U$ is constant. Hence $N_p(4,r)\le p$.

For $t=5$ the remaining classes are $U,TU,T^2U,U^2,TU^2,U^3$. First, $U^2$ lies in the span
of $1,T,T^2,U,TU$: use $U^2=2TU-T^2$ if $V^2=0$, the two-residue relation, or, in the only
other profile $(3,1,1)$, the relation $(U-c)V=0$. Next, $T^2U$, $TU^2$ and $U^3$ lie in the
span of $1,T,T^2,T^3,U,TU$. If $V^2=0$ and there are at most three residues, $U^3$ is
reduced to degree two in $U$, and $T^3-U^3=3U^2V$ (with $3$ a unit) eliminates both mixed
classes, since $T^2U=U^3+2U^2V$ and $TU^2=U^3+U^2V$. If $c$ is the only repeated residue (or
there is none, so that $V=0$), then
\[
 T^2U=cT^2+U^3-cU^2,\qquad TU^2=c^2T+U^3-c^2U,
\]
and $U^3$ is reduced by the residue polynomial when there are at most three residues, and
otherwise by $T^3-U^3=3c^2V$, which holds because then $V^2=0$ and $U^2V=c^2V$. The only
profile left is $(3,2)$, with triple residue $c$ and double residue $c'$. There $U^2$ and
$U^3$ are affine in $U$, so $TU^2$ is in the span of $T$ and $TU$. For $T^2U$, choose a
cubic $\gamma(T)$ whose derivative is the quadratic Hermite interpolant with
$\gamma'(c)=2c^2$, $\gamma''(c)=2c$ and $\gamma'(c')=2c'^2$; integration is possible since
$2$ and $3$ are units. Then $\gamma(T)-T^2U$, which equals $\gamma(T)-cT^2$ on the factor at
$c$ and $\gamma(T)-c'T^2$ on the factor at $c'$, is constant on each factor, and an affine
function of $U$ corrects these two constants. Every identity used involves only earlier
indices of no greater digit sum. Hence the largest remaining index is $p+1$ and
$N_p(5,r)\le p+1$. If $r>1/(p-3)$, then $w(p+1)-w(4)=(p-3)r-1>0$ and the index $p+1$ is
excluded by \cref{smd:lem:excluded}(a), so $N_p(5,r)\le p$.
\end{proof}

\begin{proof}[Proof of the upper bounds in \cref{smd:thm:sixroot}]
Here $p\ge7$ and the remaining classes are $T^iU^j$ with $j\ge1$, $i+j\le4$. We show that
all of them except $U$ and $TU$, and except $T^2U$ and $U^2$ in profile $(3,2,1)$, are
excluded by \cref{smd:lem:excluded}(b).

\emph{At most two residues $c,c'$.} Every $U^j$ is affine in $U$, so it suffices to treat
$T^2U$ and $T^3U$. If $c'$ has multiplicity at most two, $(U-c)(T-c')^2=0$ expresses
$T^2U$ through $T^2,TU,U,T,1$. In profile $(3,3)$ the class
$T^2U-T^3/3-(c+c')TU+cc'T$ has vanishing first and second derivatives at both residues, so it
is affine in $U$. For $T^3U$ use $(U-c)(T-c')^3=0$ with $c'$ of multiplicity at most three.
With a single residue $U$ is constant.

\emph{Exactly one repeated residue $c$, at least three residues.} Then $U^2=TU-cT+cU$, and
for $i,j\ge1$, checking on each factor,
\[
 T^iU^j=U^{i+j}+c^j\bigl(T^i-U^i\bigr).
\]
For $d=3,4$ the class $U^d$ is reduced by the residue polynomial if there are at most $d$
residues; otherwise the multiplicity of $c$ is at most $d$, and $(T-c)^d=(U-c)^d$ reduces
$U^d$ to $T^d$ and lower powers of $U$. If no residue is repeated, $U=T$.

\emph{Other profiles with $V^2=0$}, namely $(2,2,2)$ and $(2,2,1,1)$. Then $U^2=2TU-T^2$ and
$T^iU^j=U^{i+j}+iU^{i+j-1}V$, while $T^d=U^d+dU^{d-1}V$; so the mixed classes of total degree
$d$ reduce to $T^d$ and $U^d$. The class $U^d$, $d=3,4$, is reduced by the residue polynomial
when there are at most $d$ residues. In the remaining case $d=3$, profile $(2,2,1,1)$ with
double residues $c,c'$, the class $V$ is supported on two residues, so
$U^2V=(c+c')UV-cc'V$, and $T^3=U^3+3U^2V$ reduces $U^3$.

\emph{Profile $(3,2,1)$.} Let $c_3$, $c_2$ and $c_1$ be the triple, the double and the
simple residue. The six classes $1,T,T^2,U,TU,U^2$ span the same space as $1,U,U^2,V,UV,V^2$, and these are
independent: in
\[
 \mathcal R\cong\F_p[T]/(T-c_3)^3\times\F_p[T]/(T-c_2)^2\times\F_p
\]
a class $g(U)+h(U)V+k(U)V^2$ with $\deg g\le2$, $\deg h\le1$, $\deg k=0$ is determined by its values
at the three residues, its first jets at two of them and its second jet at one. So they form a
basis, which expresses $TU^2$, $U^3$, $T^2U^2$, $TU^3$ and $U^4$ through earlier indices of
digit sum at most two. Let $J$ be the functional vanishing on $1,T,T^2,U,TU$ with
$J(U^2)=1$. From $J(V)=0$, $J(UV)=-1$, $J(V^2)=1$ and the relations
$(U-c_3)(U-c_2)(U-c_1)=0$, $(U-c_3)(U-c_2)V=0$, $(U-c_3)V^2=0$, $V^3=0$ one finds
\[
 J(U^3)=c_1+c_2+c_3,\quad J(U^2V)=-(c_3+c_2),\quad J(UV^2)=c_3,\quad
 J(T^3)=c_3+c_1-2c_2 .
\]
If $c_3+c_1=2c_2$, an affine change of variable, which acts affinely on $U$ as well, gives
$(c_3,c_2,c_1)=(0,1,2)$, where $V^3=0$ and $U^2V^2=c_3^2V^2=0$ give
$J(T^4)=J(U^4)+4J(U^3V)=7-4=3$. Thus $U^2$ may be replaced by
$T^3$ or by $T^4$ in the basis, and the modified basis, of indices $0,1,2,3\text{ or }4,p,p+1$
and digit sums at most four, expresses $T^3U$.

Hence the largest remaining index is $2p$, which gives $N_p(6,r)\le2p$. If $r>1/(p-2)$, every
weight of the modified basis is smaller than $w(p+2)$ and $w(2p)$, because
$w(p+2)-w(4)=(p-2)r-1>0$, $w(2p)-w(4)=2((p-2)r-1)$ and $w(2p)-w(p+1)=(p-1)r-1>0$; by
\cref{smd:lem:excluded}(b) the indices $p+2$ and $2p$ are excluded and $N_p(6,r)\le p+1$. If
$r>1/(p-4)$ then $w(p+1)-w(5)=(p-4)r-1>0$, and if $r>1/(p-5)$ then
$w(p)-w(5)=(p-5)r-1>0$; \cref{smd:lem:excluded}(a) gives the bounds $p$ and $5$.
\end{proof}

\paragraph{Witnesses.}
For $3\le t\le6$ and $p\ge t$ put $d=p-t+1$ and
\[
 f_t(X)=X^p-\sum_{j=0}^{t-2}\binom pj(X-1)^j+p^{2t-1},
\]
a polynomial with at most $t$ terms. In $Z=X-1$ it equals
$\sum_{j=t-1}^p\binom pjZ^j+p^{2t-1}$, whose Newton polygon has vertices $(0,2t-1)$,
$(t-1,1)$, $(p,0)$ and slopes $-2$ and $-1/d$. The normalized reductions are a binomial with
$t-1$ simple nonzero roots and $Y^{t-1}(c+Y^d)$, $c\ne0$, with $d$ simple nonzero roots, as
$p\nmid(t-1)d$. Hensel's lemma gives $p$ distinct roots in the disk of depth $1/d$ about $1$.
This attains $N_p(t,r)\ge p$ for $t=3,4$ and $r\le1/(p-t+1)$, for $t=5$ and
$r\le1/(p-4)$, and for $t=6$ and $r\le1/(p-5)$. For the deepest range in each case a
monic polynomial with $t-1$ distinct integer roots $p$-adically close to $1$ suffices.

For five terms and $r\le1/(p-3)$ take
\[
 F(X)=X^{2p}+1-p^2\bigl(X^{p+1}+X^{p-1}\bigr)+\bigl(2(p^2-1)+p^5\bigr)X^p .
\]
After division by $X^p$ and the substitution $X=\exp(bY)$ the constant coefficient is $p^5$,
the odd coefficients and the quadratic one vanish, and for even $k\ge4$ the coefficient is
$2(p^k-p^2)b^k/k!$, of valuation $2+w(k)$. At $r=1/(p-3)$ the only minimal indices are $4$
and $p+1$: for even $4<k<p$ the weights increase, and for even $k>p+1$ the digit sum is even
and at least two, so $w(k)>w(p+1)$. The normalized reduction $Y^4(c+c'Y^{p-3})$,
$cc'\ne0$, gives $p-3$ simple roots of depth exactly $1/(p-3)$. At depth $3/4$ the only
minimal indices are $0$ and $4$, since $2+w(4)=5$ and
$w(k)\ge(3/4-1/(p-1))k+2/(p-1)>3$ for even $k\ge6$; this gives four further simple roots of
depth $3/4$. The two sets are disjoint, so $F$ has $p+1$ distinct roots in the disk of depth
$1/(p-3)$.

For six terms and $r=1/(p-2)$ use the exponents $0,p,2p,1,p+1,2$, of profile $(3,2,1)$ with
residues $(0,1,2)$. By the proof above the classes $1,T,T^2,T^p,T^{p+1},T^{2p}$ form a basis
of $\Z_{(p)}[T]/\prod_j(T-n_j)$. Let $L$ be the functional with values $0,0,0,0,0,1$ on this
basis; over $\Q$ the evaluations at the six exponents span the dual space, so
$L(f)=\sum_jc_jf(n_j)$ with rational $c_j$, and $P=\sum_jc_jX^{n_j}$ has moments $M_k=L(T^k)$
about the centre $1$. Then $M_0=M_1=M_2=M_p=M_{p+1}=0$, $M_{2p}=1$, and modulo $p$ the
moments $M_3,M_4,M_{p+2}$ are $J(T^3)=0$, $J(T^4)=3$ and
$J(T^2U)=J(U^3)+2J(U^2V)+J(UV^2)=1$. At $r=1/(p-2)$ the minimum coefficient valuation is
$4r$, attained exactly at $4$, $p+2$ and $2p$: the index $3$ has valuation at least
$1+3r>4r$, the index $5$ has weight $5r$, and all other indices are excluded by the proof
above. Weierstrass preparation gives $2p$ roots in the closed disk, counted with
multiplicity.

For six terms and $r=1/(p-4)$ use the exponents $0,p,2p,1,2,c$, where $c$ is the least
positive residue of $-2/3$ modulo $p$ (so $c\notin\set{0,1,2}$), and let $L$ extract the
coefficient of $T^5$ in $\Q[T]/\prod_j(T-n_j)$. In the residue algebra
\[
 U=\beta T^4-\beta(3+c)T^3,\qquad \beta=(2c)^{-1},
\]
because the right side minus $T$ is $\beta T(T-1)(T-2)(T-c)$, so the two sides agree on all
four factors. Hence $p\mid M_p$, while $M_5=1$ and $M_{p+1}\equiv\beta$ are units, and
$M_0=\dots=M_4=0$. The profile $(3,1,1,1)$ leaves only the classes $U$ and $TU$, and
$1+w(p)>w(5)$. So the minimizing indices are $5$ and $p+1$, and preparation gives $p+1$
roots.

\paragraph{Irreducible witnesses.}
In the last two constructions clear denominators. In each construction keep the leading
monomial and perturb the other coefficients, within their support, by multiples of a high
power of $p$, imposing the Eisenstein congruences at a prime different from $p$ that does not
divide the leading coefficient, by the Chinese remainder theorem. A sufficiently small
perturbation changes the analytic coefficients by terms of valuation at least the chosen
precision plus $w(k)\ge0$, so the minimizing indices at the endpoint are preserved (for
$f_t$ use its Newton polygon). The perturbed polynomial is irreducible over $\Q$, hence
separable, so its roots are distinct. For the products of linear factors the same
perturbation is applied in disjoint small disks about the original roots. Each endpoint
witness also serves at every shallower depth of its range, and a general centre is reached by
scaling the variable. This completes the proofs of \cref{smd:thm:rootcount,smd:thm:sixroot}.
\qed

The residue spans and Newton profiles are checked by exact scripts:
\smdcert{\artifact{padic-fewnomial-root-count/verify.py} with output
\artifact{padic-fewnomial-root-count/verification.txt};
\artifact{padic-six-term/verify.py} with output \artifact{padic-six-term/verification.txt}.}

\subsection{Separated complex conjugates}\label{smd:sec:arch}
{\raggedright\sources{grh_cm_prime_applications; grh_cm_four_terms_and_prime_bias; grh_sparse_singular_moduli}}

The archimedean counterpart of the disk lemma uses conjugates of very different sizes.

\begin{lemma}[Separated conjugate sizes]\label[lemma]{smd:lem:separated}
Let $z_1,\dots,z_t$ be nonzero complex numbers whose moduli $r_i=\abs{z_i}$ satisfy
$r_{i+1}\ge16r_i$. For any integers $0\le e_1<\dots<e_t$ the matrix
$(z_i^{e_j})_{i,j=1}^t$ is invertible. In particular a nonzero polynomial vanishing at all
these points has more than $t$ terms.
\end{lemma}

\begin{proof}
Multiply column $j$ by a positive number $c_j$, where $c_1=1$ and
\[
 c_{j+1}/c_j=(r_jr_{j+1})^{-(e_{j+1}-e_j)/2} .
\]
In row $i$ the ratio of successive
absolute entries is $(r_i/\sqrt{r_jr_{j+1}})^{e_{j+1}-e_j}$. It is at least $4$ up to the
diagonal ($j<i$) and at most $1/4$ after it ($j\ge i$). The sum of the off-diagonal
magnitudes in a row is at most $2\sum_{m\ge1}4^{-m}=2/3$ of its diagonal magnitude.
Strict diagonal dominance implies invertibility, and column scaling is reversible. For a
polynomial with fewer than $t$ terms, enlarge its set of exponents to $t$ exponents.
\end{proof}

When the moduli are closer, the full permutation expansion can still be compared with its
diagonal term.

\begin{lemma}[Permutation test]\label[lemma]{smd:lem:perm}
Let $z_1,\dots,z_t$ be complex numbers with moduli $0<r_1<\dots<r_t$ such that
\begin{equation}\label{smd:eq:perm}
 \sum_{\sigma\ne\mathrm{id}}\ \prod_{i=1}^tr_i^{\sigma(i)-1}<\prod_{i=1}^tr_i^{\,i-1},
\end{equation}
the sum being over the nonidentity permutations of $\set{1,\dots,t}$. Then
$\det(z_i^{e_j})\ne0$ for all integers $0\le e_1<\dots<e_t$.
\end{lemma}

\begin{proof}
Write $g_j=e_{j+1}-e_j\ge1$. Dividing the absolute value of the term of $\sigma$ in the
determinant by that of the diagonal term gives
\[
 \prod_{j=1}^{t-1}\left(\frac{\prod_{i:\,\sigma(i)>j}r_i}{\prod_{i>j}r_i}\right)^{g_j}.
\]
Each numerator has as many factors as its denominator, which is the product of the largest
radii of that cardinality, so each bracket is at most one. The quotient is therefore
largest when all gaps equal one, where it is the quotient of the corresponding terms in
\eqref{smd:eq:perm}.
\end{proof}

In applications \eqref{smd:eq:perm} is verified with rigorous enclosures
$L_i\le r_i\le U_i$ of the moduli, using the upper endpoints on the left and the lower
endpoints on the right.

\paragraph{Sizes of conjugates.}
For a primitive reduced form $[a,b,c]$ of discriminant $-D$ write $j_a$ for the associated
conjugate. The cusp estimate of \cite[\S2]{Bilu2022Trinomials} gives
\begin{equation}\label{smd:eq:cusp}
 \Bigl|\,\abs{j_a}-\exp\bigl(\pi\sqrt D/a\bigr)\Bigr|\le2079 .
\end{equation}
For two leading coefficients $a<a'$ of reduced forms, the integer inequalities
\begin{equation}\label{smd:eq:septest}
 D\ge16a^2,\qquad 9D(a'-a)^2\ge16a^2a'^2
\end{equation}
imply $\abs{j_a}>16\abs{j_{a'}}$. Indeed the exponentials $U=\exp(\pi\sqrt D/a)$ and
$V=\exp(\pi\sqrt D/a')$ satisfy $U\ge e^{4\pi}>34\cdot2079$ and
$U/V\ge e^{4\pi/3}>32$, whence $U-2079>16(V+2079)$. A chain of $t$ leading
coefficients whose consecutive pairs pass \eqref{smd:eq:septest} therefore yields $t$
conjugates to which \cref{smd:lem:separated} applies, and both inequalities persist when
$D$ increases.

\begin{lemma}[Prescribed leading coefficients]\label[lemma]{smd:lem:forms}
Let $a=t^2s$, where $t\mid f$ is odd and squarefree, $s$ is odd and coprime to $t$, and
every prime dividing $s$ splits in $K$ and does not divide $D$. If $D\ge4a^2$, there is a
primitive reduced form of discriminant $-D$ with leading coefficient $a$.
\end{lemma}

\begin{proof}
Put $D'=D/t^2$ and choose $v$ as follows. Its parity is that of $D'$, so that
$4\mid v^2+D'$. Modulo each prime power dividing $s$, it is a lift by Hensel's lemma of a
unit square root of $-D'$. For each prime $q\mid t$ it avoids the at most two roots of
$v^2+D'\equiv0\pmod q$. The Chinese remainder theorem gives a representative with
$\abs v\le ts$. Take $b=tv$ and $c=(v^2+D')/(4s)$. Then $b^2-4ac=-D$ and
$\abs b\le a\le c$, the last inequality because $D'\ge4t^2s^2$. The choices at the primes
dividing $t$ ensure $q\nmid c$, and those at the primes dividing $s$ ensure $q\nmid b$.
Thus $[a,b,c]$ is primitive and reduced, after normalizing the sign of $b$ on the boundary
of the fundamental domain, which changes neither $a$ nor the modulus of the conjugate.
\end{proof}

For odd $p$ one has $p\mid f$ exactly when $p^2\mid D$, because the odd part of $d_K$ is
squarefree; the hypotheses of the lemma are therefore congruence conditions on $D$.

\subsection{Reduction of the conjugates at a nonsplit prime}\label{smd:sec:reduction}
{\raggedright\sources{grh_cm_four_terms_and_prime_bias; grh_cm_sparse_equations; grh_sparse_singular_moduli}}

\begin{lemma}[Supersingular disks]\label[lemma]{smd:lem:reduction}
Let $p$ be a prime not dividing $f$ which does not split in $K$, and fix a place above $p$
of the ring class field $L$ of the order.
\begin{enumerate}[label=\textup{(\roman*)}]
\item Every conjugate of discriminant $\Delta$ reduces to a supersingular $j$-invariant in
characteristic $p$.
\item The absolute ramification index of $L$ at $p$ is at most two, and equals one when
$p$ is inert. Consequently two conjugates with the same nonzero reduction lie in a common
multiplicative disk of relative depth at least $1/2$, and at least $1$ when $p$ is inert.
\item The number of supersingular $j$-invariants over $\overline{\F}_p$ is $p/12+O(1)$.
For $p\equiv1\pmod3$ the invariant $0$ is ordinary, and the number of supersingular
invariants is $S_p=\floor{p/12}+\mathbf 1_{p\equiv7\ (\mathrm{mod}\ 12)}$.
\item In characteristics $7$ and $13$ the supersingular invariant is unique and equals
$6$ and $5$ respectively. In characteristic $19$ there are two, namely $18$ and $7$.
\end{enumerate}
In particular, if $p\in\set{7,13}$ does not divide $f$ and does not split in $K$, all
$h$ conjugates lie in one multiplicative disk about a unit, of relative depth at least
$1/2$, and at least $1$ when $(\Delta/p)=-1$.
\end{lemma}

\begin{proof}
The conjugates are algebraic integers and the corresponding curves have potentially good
reduction. Part (i) is Deuring's reduction theorem
\cite{Elkies2005Reduction,Belding2008Computing}. The field $L$ is unramified over $K$ away
from $f$ \cite{Cox2013Primes}, and $p$ has ramification index at most two in $K$; this
gives the first assertion of (ii). If $\alpha,\alpha'$ are conjugates with the same
nonzero reduction, then $\alpha'$ is a unit and $v_p(\alpha-\alpha')>0$ lies in
$\tfrac12\Z$, or in $\Z$ in the unramified case, so $v_p(\alpha/\alpha'-1)\ge1/2$,
respectively $\ge1$. The count in (iii) is classical
\cite{Elkies2005Reduction,Silverman2009Arithmetic}. In characteristic $p>3$ the curve
$y^2=x^3+Ax+B$ is supersingular exactly when the coefficient of $x^{p-1}$ in
$(x^3+Ax+B)^{(p-1)/2}$ vanishes \cite[Ch.~V, \S4]{Silverman2009Arithmetic}. For
$p\equiv1\pmod3$ and $j=0$ this coefficient is
$\binom{(p-1)/2}{(p-1)/3}\not\equiv0\pmod p$. In characteristic $7$ it is $3B$, giving
$j=1728\equiv6$. In characteristic $13$ it is $20A^3+15B^2$, giving $B^2=3A^3$ and
$j=1728\cdot4/(4+81)=5$; here $A\ne0$ for a nonsingular curve. In characteristic $19$
it is $B(10A^3+8B^2)$, and vanishing gives the values $18$ and $7$ for nonsingular
curves.
\end{proof}

\subsection{A finite test for unbounded exponents}\label{smd:sec:confluent}
{\raggedright\sources{grh_sparse_singular_moduli; grh_transfer_bridges_2026-11-25}}

When no single disk and no chain of separated conjugates is available, linear independence
of powers with arbitrarily large exponents can still be reduced to finitely many rank
computations over finite fields. The test below retains information when several
exponents share a residue.

\begin{theorem}[Blockwise confluent certificate; \stU]\label{smd:thm:confluent}
Let $p$ be any prime, let $\sigma,m\ge1$ be integers, and let $A$ be a finite free
commutative $\Z_{(p)}$-algebra. Suppose
\[
 x^m=a+p^\sigma D_0,\qquad p\nmid a,
\]
with $x,D_0\in A$ and $a\in\Z$. Let $\mathcal E$ be a finite set of distinct nonnegative
integers, and let $n_r$ count its elements congruent to $r$ modulo $m$, for $0\le r<m$.
Assume
\begin{equation}\label{smd:eq:block}
 n_r<p^\sigma\qquad\text{for every }r .
\end{equation}
If the columns
\[
 x^rD_0^{\,l},\qquad0\le r<m,\quad0\le l<n_r,
\]
are linearly independent in $A/pA$ over $\F_p$, then the powers $x^e$, $e\in\mathcal E$,
are linearly independent in $A\otimes\Q$ over $\Q$.
\end{theorem}

\begin{proof}
Within a residue block write the exponents as $e_i=r+mt_i$ and divide the associated
columns by $a^{t_i}$. They become $x^r(1+p^\sigma D_0/a)^{t_i}$. For
$l=0,\dots,n_r-1$ take the Newton divided difference of order $l$ at the first $l+1$
distinct integers $t_i$, multiplied by $l!/p^{\sigma l}$. These are invertible triangular
column operations over $\Q$.

Expand by the binomial theorem up to the largest $t_i$ in the block. The coefficient of
$(D_0/a)^i$ in the new column of order $l$ is
\[
 l!\,p^{\sigma(i-l)}\,[t_0,\dots,t_l]\binom Ti .
\]
It is zero for $i<l$ and one for $i=l$. For $i>l$ its valuation is at least
$\sigma(i-l)-v_p(i!/l!)$. Indeed $i!\binom Ti$ is an integer polynomial, and divided
differences of integer polynomials at integer arguments are integers: for a monomial the
divided difference of order $l$ is the complete homogeneous symmetric polynomial of the
complementary degree at the chosen arguments, and linearity gives the general case.

Since $l\le n_r-1\le p^\sigma-2$, the digit sums satisfy $S_p(l)\le\sigma(p-1)-1$ and
$S_p(i)\ge1$. Put $d=i-l\ge1$ and $C=\sigma(p-1)-1\ge0$. Legendre's formula gives
\[
 (p-1)\bigl(\sigma d-v_p(i!/l!)\bigr)=Cd-S_p(l)+S_p(i)\ge C(d-1)+1>0 .
\]
Every coefficient with $i>l$ is therefore divisible by $p$. The transformed column is
integral and reduces to $x^r(D_0/a)^l$ modulo $p$, and the factor $a^{-l}$ is a nonzero
scalar modulo $p$.

A rational dependence between integral columns in a finite free $\Z_{(p)}$-module can be
scaled to have integral coefficients with at least one unit. Its reduction would
contradict the assumed independence. Reversing the invertible column operations gives the
result.
\end{proof}

The proof uses finite binomial expansions only; no $p$-adic exponential or logarithm is
required, so $p=2$ causes no convergence problem. The bound \eqref{smd:eq:block} is on each
residue block, not on the total support. In the applications of
\cref{smd:sec:quintic,smd:sec:sextics} one has $\sigma=1$, $A=\Z_{(p)}[X]/(H_\Delta)$, $x$
the class of $X$, and at most five exponents, so that \eqref{smd:eq:block} holds for every
$p\ge7$. The congruence $x^m\equiv a\pmod{pA}$ is certified by reducing $x^m$ modulo
$(H_\Delta,p^2)$: the constant coefficient is a $p$-adic unit, every other coefficient is
divisible by $p$, and taking $a$ equal to the constant coefficient determines
$D_0\bmod p$.

\begin{proposition}[Sharpness of the strict threshold]\label[proposition]{smd:prop:sharp}
The strict inequality in \eqref{smd:eq:block} cannot be replaced by a weak inequality for
all finite free algebras. Put $k=p^\sigma$ and
\[
 A=\Z_{(p)}[D_0]/(D_0^k+D_0^{k-1}-1),\qquad x=1+kD_0,\quad m=1 .
\]
The $k$ confluent columns $1,D_0,\dots,D_0^{k-1}$ form a basis of $A/pA$, but
\[
 F(X)=(X-1)^k+k(X-1)^{k-1}-k^k
\]
vanishes at $x$ and has exactly $k$ nonzero coefficients.
\end{proposition}

\begin{proof}
One has $F(x)=k^k(D_0^k+D_0^{k-1}-1)=0$. The coefficient of $X^{k-1}$ in $F$ is
$-k+k=0$, and the coefficients at the exponents $0,1,\dots,k-2,k$ are nonzero. Hence
the $k$ powers $x^e$, $e\in\set{0,\dots,k-2,k}$, are dependent, although they lie in one
residue block with $n_0=k=p^\sigma$.
\end{proof}

At $p=3$, $\sigma=1$ the example becomes the irreducible relation $x^3-3x-25=0$ after
$D_0^3+D_0^2-1=0$. The proposition concerns the generality of finite free algebras; it
does not assert that $A=\Z_{(p)}[x]$, and that equality need not hold.

\subsection{Equations with at most five terms: the classification}\label{smd:sec:classification}
{\raggedright\sources{grh_sparse_singular_moduli; grh_cm_sparse_equations; grh_cm_four_terms_and_prime_bias; grh_cm_prime_applications}}

\begin{theorem}[Five-term classification; \stU]\label{smd:thm:five}
For every imaginary quadratic order discriminant $\Delta$,
\[
 s(\Delta)\le5\quad\Longleftrightarrow\quad h(\Delta)\le4 .
\]
Moreover $s(\Delta)=h(\Delta)+1$ whenever $h(\Delta)\le5$, with the single exception
$\Delta=-3$, where the singular modulus is $0$ and $s=1$. In particular every quartic
singular modulus requires exactly five terms and every quintic singular modulus exactly
six. These assertions are independent of RH and GRH.
\end{theorem}

The content of the reverse implication is that for every singular modulus of degree at
least five, any five distinct nonnegative powers are linearly independent over $\Q$.

\begin{corollary}[Three and four terms; \stU]\label[corollary]{smd:cor:riffaut}
A singular modulus is a root of a nonzero rational polynomial with at most three terms if
and only if its degree is at most two, and of one with at most four terms if and only if
its degree is at most three. In particular Riffaut's conjecture holds without any
hypothesis on zeros.
\end{corollary}

\begin{proof}
For $h\le3$ use \eqref{smd:eq:trivial}. For $h=3,4$ the theorem gives $s=4,5$, and for
$h\ge5$ it gives $s\ge6$.
\end{proof}

\begin{remark}[The trinomial case]\label[remark]{smd:rem:riffaut}
The original route to the trinomial statement assumed Dirichlet GRH: by
\cite[Theorems~1.3, 1.4 and Corollary~5.2]{Bilu2022Trinomials} a discriminant admitting a
trinomial is fundamental, satisfies $D>10^{11}$, and has $p>4\sqrt D/\log D$ for every
prime $p$ splitting in $K$, while their Proposition~9.2 produces a smaller split prime
under the hypothesis supplied by \cref{src:thm:dgrh}. The unconditional reductions just
quoted can also be completed at a single prime, without the finite computations behind
\cref{smd:thm:five}. Indeed $7$ is then inert or ramified in $K$, and $f=1$. By
\cref{smd:lem:reduction} every conjugate lies in the disk $v_7(x/6-1)\ge1/2$. Three
distinct conjugates in this disk contradict \cref{smd:cor:threshold} with $t=3$ and
$\kappa(7,3)=1/5$. A binomial is covered by allowing a zero middle coefficient, or by
squaring it.
\end{remark}

The proof of \cref{smd:thm:five} occupies
\cref{smd:sec:small,smd:sec:quintic,smd:sec:global,smd:sec:sextics}. The forward
implication is \eqref{smd:eq:trivial}. Degrees up to five are settled by a complete census
(\cref{smd:prop:small,smd:thm:131}). For $h\ge5$, a combination of local and archimedean
certificates covers all but $71$ discriminants (\cref{smd:prop:global}), and these are
treated individually (\cref{smd:prop:residual,smd:prop:sextics}).

\subsection{Degrees at most five}\label{smd:sec:small}
{\raggedright\sources{grh_sparse_singular_moduli; grh_cm_sparse_equations; grh_cm_four_terms_and_prime_bias}}

\paragraph{Certified class polynomials.}
For a given discriminant, $H_\Delta$ is computed by enumerating the primitive reduced
forms and evaluating
\[
 j(\tau)=\frac{E_4(q)^3}{q\prod_{n\ge1}(1-q^n)^{24}},\qquad
 E_4(q)=1+240\sum_{n\ge1}\sigma_3(n)q^n,\quad q=e^{2\pi i\tau},
\]
with outward-rounded dyadic intervals. Reducedness gives $\abs q<1/200$. The $164$-term
truncation of $E_4$ has error at most $960\cdot165^3/200^{165}$, since
$\sigma_3(n)<2n^3$. For the omitted factors of the product the sum of $\abs q^n$ is at
most $2/200^{165}$, which gives an absolute full-product error below $8/200^{165}$. Each
coefficient enclosure of $\prod(X-j(\tau))$ contains exactly one integer, which determines
the class polynomial by integrality.
\smdcert{class polynomials
\artifact{grh_sparse_singular_moduli/cm/class_polynomials.json} and
\artifact{grh_sparse_singular_moduli/checks/five_small_witnesses.json}, with the verifier
\artifact{grh_sparse_singular_moduli/cm/verify_class_polynomials.py}; independent
reconstructions
\artifact{grh_cm_four_terms_and_prime_bias/checks/class_polynomials.json} and
\artifact{grh_cm_sparse_equations/certificates/cm-all-polys.json}.}

\paragraph{The census of orders.}
The order class number formula is
\begin{equation}\label{smd:eq:classnumber}
 h(\Delta)=\frac{h(d_K)\,f}{[\OO_K^\times:\OO_f^\times]}
 \prod_{p\mid f}\Bigl(1-\frac{\chi_{d_K}(p)}p\Bigr),
\end{equation}
where $\OO_f$ is the order of conductor $f$ and $\chi_{d_K}=(d_K/\cdot)$
\cite[Theorem~7.24]{Cox2013Primes}. The unit index is at most three, so
$h(\Delta)\ge h(d_K)\varphi(f)/3$, and $h(d_K)$ divides $h(\Delta)$ by surjectivity of the
map of class groups. On prime powers one has $\varphi(f)^2\ge f/2$, hence
$f\le2\varphi(f)^2\le18h^2$. Watkins's table \cite[Table~4]{Watkins2004Class} bounds the
fundamental discriminants: $\abs{d_K}\le163,427,907,1555,2683$ for
$h(d_K)=1,2,3,4,5$. Enumerating the fundamental discriminants, the conductors
($f\le162$, $288$, $450$ for $h=3,4,5$) and the reduced forms gives exactly $25$ cubic,
$84$ quartic and $29$ quintic orders. The quintic orders have
\[
\begin{aligned}
 D={}&47, 79, 103, 127, 131, 179, 188, 227, 316, 347, 412, 443, 508, 523, 571, 619, 683,
 691, 739,\\
 &787, 947, 1051, 1123, 1723, 1747, 1867, 2203, 2347, 2683,
\end{aligned}
\]
and the quartic orders are listed in the proof below.

\paragraph{Disk certificates.}
A \emph{disk certificate for $t$ terms} at $\Delta$ is a quadruple $(p,v,u,\lambda)$,
with $p$ prime, $v\ge0$ an integer, $u$ an integer prime to $p$ and
$\lambda>1/(p-1)$ rational, such that the least minimizing index of
\cref{smd:lem:newton} for $H_\Delta(p^v(u+Z))$ at $\lambda$ is at least $t$, and
\[
 w_{p,\lambda}(n)>\max_{i<t}w_{p,\lambda}(i)\qquad\text{for all }n\ge t .
\]
The first condition places at least $t$ conjugates in the multiplicative disk of relative
depth $\lambda$ about $p^vu$, and the second says $B_p(t,\lambda)<t$. By
\cref{smd:lem:disk} such a certificate excludes all rational relations with at most $t$
terms, of arbitrary degree, so that $s(\Delta)\ge t+1$. The second condition is a finite
check: one uses exact factorial valuations up to the point where
$n\lambda-(n-1)/(p-1)$ exceeds the required maximum, and this lower bound for
$w_{p,\lambda}(n)$ increases thereafter. \Cref{smd:ex:hand} gives two instances.

\begin{proposition}[Degrees up to five; \stU]\label[proposition]{smd:prop:small}
Let $h(\Delta)\le5$ and $\Delta\notin\set{-3,-131}$. Then $s(\Delta)=h(\Delta)+1$.
\end{proposition}

\begin{proof}
By \eqref{smd:eq:trivial} only lower bounds are needed.

\emph{Degree one.} A nonzero rational number is not a root of a monomial, so $s=2$.

\emph{Degree two.} Here $D\ge15$, and the conjugate of the principal form has larger
absolute value than the other conjugate, since
$e^{\pi\sqrt D}-2079>e^{\pi\sqrt D/2}+2079$ by \eqref{smd:eq:cusp}. A binomial has all its
nonzero complex roots on one circle, so $s=3$.

\emph{Degree three.} This is \cite[Theorem~6.6]{Bilu2022Trinomials}, which is independent
of GRH. Independently, each of the $25$ cubic orders has a disk certificate for three terms
at a prime $p\in\set{2,3,7,13,17}$.

\emph{Degree four.} For the $50$ orders with
\[
\begin{aligned}
 D={}&39, 55, 56, 63, 68, 80, 128, 136, 144, 155, 156, 171, 184, 196, 203, 208, 219, 220,
 252, 256,\\
 &259, 275, 291, 292, 323, 328, 355, 363, 387, 388, 400, 475, 507, 568, 592, 603, 667,
 723,\\
 &763, 772, 955, 1003, 1027, 1227, 1243, 1387, 1411, 1467, 1507, 1555,
\end{aligned}
\]
the certificate file listed below contains a disk certificate for four terms. For the
remaining $34$ quartic orders,
\[
\begin{aligned}
 D={}&84, 96, 120, 132, 160, 168, 180, 192, 195, 228, 240, 280, 288, 312, 315, 340, 352,
 372, 408,\\
 &435, 448, 483, 520, 532, 555, 595, 627, 708, 715, 760, 795, 928, 1012, 1435,
\end{aligned}
\]
certified enclosures of the four moduli satisfy \eqref{smd:eq:perm}, and
\cref{smd:lem:perm} applies.

\emph{Degree five.} For $28$ of the $29$ quintic orders the same file contains a disk
certificate for five terms. The witness primes are distributed as follows.
\begin{center}
\begin{tabular}{lrrrrrr}
\toprule
Witness prime&2&3&7&13&17&Total\\
\midrule
Four-term exclusions (50 quartic and 29 quintic orders)&1&48&23&5&2&79\\
Five-term exclusions by a single disk (quintic orders)&0&20&5&3&0&28\\
\bottomrule
\end{tabular}
\end{center}
The remaining quintic order is $\Delta=-131$, excluded from the statement.
\end{proof}

\smdcert{cubic orders
\artifact{grh_sparse_singular_moduli/cm/verify_all_cubics.py}; disk certificates
\artifact{grh_sparse_singular_moduli/cm/disk_witnesses.csv}, checked by
\artifact{grh_sparse_singular_moduli/cm/verify_disks.py}; permutation tests for the $34$
quartic orders \artifact{grh_sparse_singular_moduli/cm/verify_all_quartics.py}.}

Independent sets of disk certificates for the same $79$ orders exist. One of them uses
only the eight triples $(p,v,u)=(2,3,5)$, $(2,6,3)$, $(2,6,5)$, $(3,3,1)$, $(3,3,2)$,
$(7,0,6)$, $(13,0,5)$, $(17,0,8)$ with the thresholds of \cref{smd:cor:threshold}, and is
extended to a disk for each quartic order.
\smdcert{\artifact{grh_cm_four_terms_and_prime_bias/checks/disks.csv};
\artifact{grh_cm_sparse_equations/certificates/cm-clusters.csv} (eight triples);
\artifact{grh_cm_sparse_equations/certificates/cm-all-clusters.csv} (all quartic and
quintic orders).}

At $\Delta=-131$ the reduction of $H_{-131}$ modulo $17$ is $X(X-8)^4$, and the
$17$-adic valuations of the coefficients of $H_{-131}(8+Y)$, constant term first, are
$(4,3,2,1,0,0)$. At $\lambda=1/14$ the least minimizing index is $4$, so four conjugates
lie at depth greater than $\kappa(17,4)=1/14$ about $8$ and $s(-131)\ge5$; the fifth
conjugate has valuation $3$ and lies outside every such disk. A five-term relation
therefore requires a different argument.

\subsection{The last quintic order}\label{smd:sec:quintic}
{\raggedright\sources{grh_sparse_singular_moduli; grh_transfer_bridges_2026-11-25}}

Let $H=H_\Delta$ have degree $h$ and let $x$ be a root. Suppose that five powers of $x$
are linearly dependent over $\Q$. Since $x\ne0$, dividing by the lowest power gives
exponents $0=e_0<e_1<e_2<e_3<e_4$. For a modulus $N$, the \emph{pattern} of the relation
modulo $N$ is the sorted quadruple of the residues of $e_1,\dots,e_4$, in which
repetitions and zeros are allowed. Let $(p,m)$ be a pair with $m\mid N$, $p\ge7$ and
$x^m\equiv a\pmod{pA}$ for a rational integer $a$ prime to $p$, in the notation following
\cref{smd:thm:confluent}. The pattern modulo $N$ determines the numbers $n_r$ of
\cref{smd:thm:confluent} for the five exponents, the exponent $0$ included, and by that
theorem the five confluent columns must be linearly dependent in $A/pA$. The exponent $0$
supplies the column $1$; in the basis $1,x,\dots,x^{h-1}$ the test therefore amounts to
the vanishing of all $4\times4$ minors of the matrix of the other four columns with the
constant coordinate deleted. A pattern passing this test for every pair considered at a
given stage is said to \emph{survive}.

If $n\mid N$, every pattern modulo $N$ is obtained from its reduction modulo $n$ by adding
multiples of $n$ to the four residues. For equal old residues it suffices to take
nondecreasing lifts among those positions, and sorting the resulting quadruple gives every
possible lift exactly once. A genuine exponent set gives a surviving pattern at every
stage of such a chain of moduli. If no pattern survives at the last stage, any five
distinct powers of $x$ are independent, and a relation on fewer powers is excluded by
adjoining further powers.

\begin{theorem}[\stU]\label{smd:thm:131}
The least support for discriminant $-131$ is six: $s(-131)=6$.
\end{theorem}

\begin{proof}
Apply the sieve just described to the certified polynomial $H_{-131}$, of degree five,
with the following stages. The modulus at every stage is divisible by each period used
there.
\begin{center}
\begin{tabular}{rlrr}
\toprule
$N$ & Pairs $(p,m)$ tested at this stage & Patterns tested & Surviving\\
\midrule
200    & $(601,200)$                    &68,685,050 &119,264\\
1,400  & $(1051,350),\ (4201,1400)$      &279,218,674&196\\
4,200  & $(6301,1050)$                  &15,876     &5,876\\
12,600 & $(757,252)$                    &475,956    &214,120\\
25,200 & $(337,112)$                    &3,425,920  &0\\
\bottomrule
\end{tabular}
\end{center}
The last count proves that any five powers of a root are independent, and $H_{-131}$
itself has at most six terms. The polynomial reductions and determinant checks are
reproduced by two implementations, the second with a separately written basis recurrence,
cofactor computation and lift enumerator; both reproduce all five rows.
\end{proof}

\smdcert{\artifact{grh_sparse_singular_moduli/checks/verify_last_quintic.py};
\artifact{grh_sparse_singular_moduli/checks/independent_quintic.py}; output
\artifact{grh_sparse_singular_moduli/checks/last_quintic.txt}.}

\Cref{smd:prop:small,smd:thm:131} give $s=h+1$ for every order with $h\le5$ other than
$\Delta=-3$, and in particular $s\ge6$ for $h=5$.

\subsection{Reduction of the global problem to 71 discriminants}\label{smd:sec:global}
{\raggedright\sources{grh_sparse_singular_moduli; grh_cm_sparse_equations; grh_cm_four_terms_and_prime_bias}}

\begin{proposition}[Global reduction; \stU]\label[proposition]{smd:prop:global}
Let $h(\Delta)\ge5$. Then $s(\Delta)\ge6$, except possibly when $D$ is one of the
following $71$ integers, all with $h\ge6$ and $87\le D\le4563$:
\[
\begin{aligned}
 &87, 152, 244, 248, 251, 264, 272, 276, 300, 339, 348, 360, 376, 432, 451, 467, 472, 563,
 588,\\
 &608, 612, 628, 640, 664, 675, 699, 703, 712, 724, 755, 768, 784, 859, 867, 972, 979,
 1000,\\
 &1067, 1083, 1147, 1195, 1252, 1291, 1312, 1323, 1368, 1459, 1483, 1587, 1648, 1699,
 1728,\\
 &1732, 1795, 1819, 1875, 1972, 2107, 2187, 2323, 2467, 2523, 2539, 2707, 2812, 3043,
 3643,\\
 &3667, 3675, 3883, 4563 .
\end{aligned}
\]
\end{proposition}

\begin{proof}
\emph{Local certificates.} Let $p\in\set{7,13}$ with $p\nmid f$ and $p$ nonsplit in $K$.
By \cref{smd:lem:reduction} all $h\ge5$ conjugates lie in one multiplicative disk of
relative depth at least $1/2$ about a unit. In \cref{smd:lem:disk} take $r=1/2$ and
$t=5$: the first five weights $w_{p,1/2}(i)=i/2$ are at most $2$, for $n=5,6$ one has
$n/2-v_p(n!)>2$, and for $n\ge7$ one has $n/2-v_p(n!)\ge n/2-(n-1)/6>2$. Hence
$B_p(5,1/2)\le4$ and $s(\Delta)\ge6$. If $7\cdot13\mid f$, the prime $19$ can be used in
the same way when $19\nmid f$ and $19$ is nonsplit: \eqref{smd:eq:classnumber} gives
$h\ge\varphi(f)/3\ge\varphi(91)/3=24$, there are only two supersingular invariants in
characteristic $19$, both nonzero, so one disk of depth at least $1/2$ holds at least
twelve conjugates, and the weight computation is unchanged.

\emph{Archimedean certificates.} The following cases are disjoint and exhaustive. In each
row, if a local certificate above applies at a prime available in that row, the case is
settled. Otherwise every such prime splits, the displayed five leading coefficients exist
by \cref{smd:lem:forms} as soon as $D\ge4a_{\max}^2$, and their consecutive pairs pass the
separation tests \eqref{smd:eq:septest} at and above the indicated bound. The last row
needs no splitting assumption. \Cref{smd:lem:separated} then gives $s(\Delta)\ge6$.
\begin{center}
\begin{tabular}{llr}
\toprule
Conductor condition & Five leading coefficients & $D$ sufficient\\
\midrule
$7\nmid f$ and $13\nmid f$ & $1,7,13,49,91$ & 38,416\\
$7\mid f$, $13\nmid f$ & $1,13,49,169,637$ & 1,623,076\\
$13\mid f$, $7\nmid f$ & $1,7,49,169,343$ & 470,596\\
$7\cdot13\mid f$, $19\nmid f$ & $1,19,49,169,361$ & 521,284\\
$7\cdot13\cdot19\mid f$ & $1,49,169,361,8281$ & 274,299,844\\
\bottomrule
\end{tabular}
\end{center}

\emph{The finite ranges.} For odd $p$ one has $p\mid f$ exactly when $p^2\mid D$.
Enumerating the integers $D\equiv0,3\pmod4$ below the bound of each row which satisfy its
divisibility conditions gives the following counts. The leading-coefficient test searches
integers $a\le251$ for a chain of five passing \eqref{smd:eq:septest}, with a directly
verified primitive reduced form at each $a$.
\begin{center}
\begin{tabular}{rrrrr}
\toprule
Row & Candidates & Local & Five leading coefficients & Remaining\\
\midrule
1 &18,705&14,908&3,667&130\\
2 &16,464& 8,820&7,634&10\\
3 & 1,364&   768&  594&2\\
4 &    30&    16&   14&0\\
5 &    45&     0&   45&0\\
\midrule
Total &36,608&24,512&11,954&142\\
\bottomrule
\end{tabular}
\end{center}
Counting all primitive reduced forms for the $142$ remaining values gives $21$ with
$h<5$ and four with $h=5$, the latter covered by \cref{smd:sec:small,smd:sec:quintic}. For
$46$ more there is a prime $p\in\set{7,13,19,31,37,43}$ with $(\Delta/p)=-1$ and
$h>4S_p$. These primes are $1\pmod3$, so by \cref{smd:lem:reduction} the conjugates
reduce to at most $S_p$ nonzero residues, one disk of relative depth at least $1$ holds at
least five of them, and \cref{smd:cor:threshold} with $\kappa(p,5)=1/(p-4)<1$ gives
$s(\Delta)\ge6$. The remaining $71$ discriminants are those listed.
\end{proof}

\smdcert{the complete integer sieve and its final list
\artifact{grh_sparse_singular_moduli/checks/global_five.py}, with output
\artifact{grh_sparse_singular_moduli/checks/global_five.txt}.}

\subsection{The residual discriminants}\label{smd:sec:sextics}
{\raggedright\sources{grh_sparse_singular_moduli; grh_cm_sparse_equations; grh_transfer_bridges_2026-11-25}}

\begin{proposition}[\stU]\label[proposition]{smd:prop:residual}
For $67$ of the $71$ discriminants of \cref{smd:prop:global}, namely all except
$D=87,348,432,675$, one has $s(\Delta)\ge6$.
\end{proposition}

\begin{proof}
For $49$ discriminants, five conjugates are selected whose moduli $r_1<\dots<r_5$,
enclosed by the dyadic evaluation of $j$ described in \cref{smd:sec:small}, satisfy
\eqref{smd:eq:perm} with $t=5$: the lower bound for the diagonal magnitude exceeds the sum
of the upper bounds for the $119$ other terms. Here successive radii may be closer than the
factor $16$ of \cref{smd:lem:separated}, and \cref{smd:lem:perm} gives independence for
arbitrary exponents. For $18$ further discriminants a disk certificate for five terms (\cref{smd:sec:small})
excludes five terms in a single disk. For
these $18$ and for the four remaining cases the class polynomial is derived from interval
enclosures and integrality. The witness file records the forms, coefficients and local
parameters, and its verifier checks that its inputs match the $71$ discriminants of the
global reduction.
\end{proof}

\smdcert{\artifact{grh_sparse_singular_moduli/checks/five_small_witnesses.json};
\artifact{grh_sparse_singular_moduli/checks/verify_five_small.py}.}

The four remaining orders have class number six.

\begin{proposition}[The four residual sextic orders; \stU]\label[proposition]{smd:prop:sextics}
For $D=87,348,432,675$ any five distinct nonnegative powers of a singular modulus of
discriminant $-D$ are linearly independent over $\Q$.
\end{proposition}

\begin{proof}
Apply the sieve of \cref{smd:sec:quintic} to the degree-six class polynomials. In the
basis $1,x,\dots,x^5$, omitting the constant coordinate leaves a $5\times4$ matrix, and a
pattern survives exactly when every $4\times4$ minor vanishes at every prime tested at
the current stage. The period $m$ in each pair $(p,m)$ is certified by reduction modulo
$(H_{-D},p^2)$. Each step from a modulus $n$ to a multiple $N$ enumerates all lifts of the
surviving patterns.
\begin{center}
\begin{tabular}{rrrlrr}
\toprule
$D$ & $n$ & $N$ & Pairs $(p,m)$ & Patterns tested & Surviving\\
\midrule
87 & 1 & 102 & $(103,102)$ &4,780,230&13,125\\
 &102 &1,530 & $(1531,1530)$ &653,917,425&625\\
 &1,530&7,650& $(151,150)$ &390,625&915\\
 &18&486& $(487,486),\ (1459,486)$ &7,971,615&0\\
348&1&486& $(487,486),\ (1459,486)$ &2,353,329,126&0\\
432&1&36& $(109,36)$ &82,251&106\\
 &36&144& $(433,144)$ &26,212&0\\
675&1&40& $(241,40)$ &123,410&1\\
 &40&280& $(13,56)$ &2,401&0\\
\bottomrule
\end{tabular}
\end{center}
For $D=87$ the $915$ survivors modulo $7650$ reduce to exactly $15$ patterns modulo
$18=\gcd(7650,486)$, and the last step for $87$ tests every lift of those $15$ patterns to
$486$. Any set of integer exponents would give compatible residues on both sides, so the
zero final count excludes every such set. For the other discriminants the successive
moduli divide one another. A second implementation treats $D=87,348$ by eliminating two
prefix columns and grouping suffix columns by their projective classes in the quotient:
for $D=87$ the $196{,}587$ simultaneous failures modulo $486$, the $13{,}125$ failures
modulo $102$ and the $19{,}669$ failures modulo $150$ restrict the lifts to $979{,}524$
candidates at modulus $1530$, all eliminated by the rank test at $1531$; for $D=348$ it
finds no survivor modulo $486$.
\end{proof}

\smdcert{\artifact{grh_sparse_singular_moduli/checks/verify_five_sieves.py}, with
enumerator \artifact{grh_sparse_singular_moduli/checks/five_sieves.cpp} and output
\artifact{grh_sparse_singular_moduli/checks/five_sieves.txt}; second implementation
\artifact{grh_sparse_singular_moduli/checks/fast_sieve/verify.py}.}

\begin{proof}[Proof of \cref{smd:thm:five}]
If $h\le4$ then $s\le5$ by \eqref{smd:eq:trivial}. If $h=5$ then $s=6$ by
\cref{smd:prop:small,smd:thm:131}. If $h\ge6$ then $s\ge6$ by
\cref{smd:prop:global,smd:prop:residual,smd:prop:sextics}. The exact values for $h\le5$
are \cref{smd:prop:small,smd:thm:131}.
\end{proof}

\begin{remark}[A second global sieve]\label[remark]{smd:rem:secondsieve}
An independent reduction reaches the same four sextic orders. With the cutoff
$A=4\cdot2401^2=23{,}059{,}204$ it tries the primes $7,13,19,31$ in order until one does
not divide $f$. If that prime is nonsplit, a disk test as in \cref{smd:prop:global}
applies (at $19$ and $31$ one has $91\mid f$, hence $h\ge24>4S_p$). If it splits, or if
all four primes divide $f$, the chains of leading coefficients $1,7,49,343,2401$;
$1,13,49,169,2197$; $1,19,49,169,361$; $1,31,49,169,361$; and $1,49,169,361,961$ are
used. Of the $11{,}529{,}601$ integers $3\le D<A$ with $D\equiv0,3\pmod4$, the sieve
discharges $9{,}109{,}793$ at a nonsplit prime $7$ or $13$ away from the conductor,
$2{,}419{,}666$ by five separated moduli, $25$ because $h<6$, and $51$ at a larger
nonsplit prime with $h>4S_p$; $66$ class polynomials with $D\le4563$ and degrees $6$ to
$12$ remain, of which $62$ are eliminated by a single five-root disk.
\smdcert{\artifact{grh_cm_sparse_equations/five/five_sieve.cpp};
\artifact{grh_cm_sparse_equations/five/counts.txt};
\artifact{grh_cm_sparse_equations/five/residual.txt};
\artifact{grh_cm_sparse_equations/certificates/cm-extended-clusters.csv}.}
\end{remark}

\begin{remark}[Limits of the single-disk test]\label[remark]{smd:rem:disklimit}
The residue sieve cannot be replaced by another choice of prime in
\cref{smd:cor:threshold}. For each of the five orders
$\Delta=-87,-131,-348,-432,-675$, no five conjugates lie in a common multiplicative disk
of depth strictly greater than $\kappa(p,5)$ at any rational prime $p$, even with an
algebraic centre. Five such roots would give $20$ ordered pairs with
$v_p(\alpha_j/\alpha_i-1)>\kappa(p,5)$. Such a pair implies $p\mid\disc(H_\Delta)$, since
the $\alpha_i$ are integral and $v_p(\alpha_j-\alpha_i)>0$. The roots of
\[
 Y^{-h}\Res_X\bigl(H_\Delta(X),H_\Delta((1+Y)X)\bigr)
\]
are the ratios $\alpha_j/\alpha_i-1$ with $i\ne j$. Factoring $\abs{\disc(H_\Delta)}$ and
applying the least-minimizing-index count to this resultant at every prime divisor shows
that the maximum number of such pairs is $14,12,14,12,14$ for the five discriminants in
the order listed.
\smdcert{\artifact{grh_cm_sparse_equations/certificates/cm_disk_limit.py}.}
\end{remark}

\subsection{Residual orders of degrees six and seven with full support}\label{smd:sec:seven}
{\raggedright\sources{singular-modulus-support; singular_moduli_seven_terms_2026-11-26; grh_cm_sparse_equations; grh_cm_prime_applications}}

For fourteen of the residual discriminants of \cref{smd:prop:global} a single disk holds all
conjugates, deep enough to exclude relations with as many terms as the degree.

\begin{theorem}[Full support for fourteen residual orders; \stU]\label{smd:thm:seven}
One has $s(\Delta)=h(\Delta)+1$ for $\Delta=-D$ in each of the following cases:
\[
\begin{aligned}
 h&=6,\ s=7: && D\in\set{152,244,300,451,472,588,628,867,1083,1147},\\
 h&=7,\ s=8: && D\in\set{859,1483,2467,2707}.
\end{aligned}
\]
In particular any six distinct integer powers of a singular modulus of discriminant $-152$
are linearly independent over $\Q$; negative exponents are allowed, after multiplication by a
common power.
\end{theorem}

\begin{proof}
For each row of the following table put $A=p^vu$ and $B=p^{v+m}$, with $p\nmid u$ and
$h=h(-D)$, and expand $p^{-vh}H_{-D}\bigl(p^v(u+Z)\bigr)=\sum_{i=0}^h\beta_iZ^i$. The last
column lists $v_p(\beta_0),\dots,v_p(\beta_h)$; in every row $\beta_h=1$ and
$v_p(\beta_i)\ge m(h-i)$.
\begin{center}
\small
\begin{tabular}{rrrrrrl}
\toprule
$D$ & $h$ & $p$ & $v$ & $u$ & $m$ & $v_p(\beta_0),\dots,v_p(\beta_h)$\\
\midrule
152  &6&2&6&1&2&12, 11, 8, 8, 4, 3, 0\\
244  &6&2&6&3&2&12, 13, 8, 10, 4, 5, 0\\
300  &6&2&4&3&2&12, 11, 8, 8, 4, 3, 0\\
451  &6&3&3&1&2&12, 12, 9, 6, 8, 3, 0\\
472  &6&2&6&1&2&12, 11, 8, 8, 4, 3, 0\\
588  &6&2&4&3&2&12, 11, 8, 9, 4, 3, 0\\
628  &6&2&6&3&2&12, 13, 8, 10, 4, 5, 0\\
867  &6&3&3&2&2&14, 12, 9, 7, 4, 5, 0\\
1083 &6&17&0&8&1&6, 5, 4, 4, 2, 1, 0\\
1147 &6&3&3&1&2&12, 11, 9, 6, 5, 3, 0\\
859  &7&3&3&1&1&14, 12, 11, 8, 6, 5, 2, 0\\
1483 &7&3&3&1&1&14, 12, 11, 8, 6, 5, 2, 0\\
2467 &7&3&3&1&1&14, 12, 11, 8, 6, 5, 2, 0\\
2707 &7&3&3&1&1&14, 12, 11, 8, 6, 5, 2, 0\\
\bottomrule
\end{tabular}
\end{center}
Hence $H_{-D}(A+BY)/B^h$ is monic with integral coefficients, and all its roots are
integral: if $v_p(Y)<0$, the leading monomial has strictly smaller valuation than every
other term. Thus every conjugate $\alpha$ satisfies $v_p(\alpha/A-1)\ge m$. For $D=152$ the
transformed polynomial is
\[
\begin{aligned}
 \frac{H_{-152}(64+256Y)}{256^6}={}&Y^6-258774476247186Y^5+259766390712739705665Y^4\\
 &+408822444341150448859180Y^3+648287850356451166578247815Y^2\\
 &-931493942482171555920209906Y+1404883983219431921036769431 .
\end{aligned}
\]

Apply \cref{smd:lem:disk} with $r=m>1/(p-1)$ and $t=h$. For $(p,m)=(2,2)$ one has
$w(k)=2k-v_2(k!)=k+S_2(k)$, whose maximum for $k<6$ is $w(5)=7$, while $w(6)=8$ and
$w(k)\ge k+1\ge8$ for $k\ge7$. For $(p,m)=(3,2)$, $w(k)=(3k+S_3(k))/2$ has maximum $9$ for
$k<6$ and exceeds $9$ for $k\ge6$. For $(p,m)=(17,1)$ the maximum for $k<6$ is $5$, and
$w(k)\ge k-(k-1)/16>5$ for $k\ge6$. For $(p,m)=(3,1)$ and $h=7$, $w(k)=(k+S_3(k))/2$ has
maximum $4$ for $k<7$, $w(7)=5$, and $w(k)\ge(k+1)/2>4$ for $k\ge8$. In every row
$B_p(h,m)<h$, so a nonzero polynomial with at most $h$ terms has fewer than $h$ roots in the
disk. A rational relation with at most $h$ terms at one conjugate vanishes at all $h$
conjugates, because $H_{-D}$ is the minimal polynomial; this is impossible. The class
polynomial gives the opposite inequality \eqref{smd:eq:trivial}.
\end{proof}

\smdcert{class polynomials from
\artifact{grh_sparse_singular_moduli/checks/five_small_witnesses.json}, reproduced in
\artifact{singular-modulus-support/polynomials.json} and recomputed from the
$q$-expansion of $j$ with interval arithmetic by
\artifact{singular-modulus-support/regenerate.py} (output
\artifact{singular-modulus-support/regeneration.txt}); the table is checked by binomial expansion and
by Horner substitution in \artifact{singular-modulus-support/verify.py} (output
\artifact{singular-modulus-support/verification.txt}) and, for the ten sextic rows, in
\artifact{singular_moduli_seven_terms_2026-11-26/verify.py} with
\artifact{singular_moduli_seven_terms_2026-11-26/certificates.json}.}

\begin{example}[A sextic order requiring seven terms]\label[example]{smd:ex:4075}
For $\Delta=-4075=-5^2\cdot163$ the class number is six, and $7$ is inert in $K$ and
prime to the conductor. By \cref{smd:lem:reduction} all six conjugates lie in one
multiplicative disk of relative depth at least $1$, which exceeds
$\kappa(7,6)=1/2$. Hence $s(-4075)=7$ by \cref{smd:cor:threshold}.
\end{example}

For $p\ge7$, \cref{smd:thm:sixroot} shows that the depth of a disk alone excludes six-term
relations at six of its points exactly when the depth exceeds $1/(p-5)$; the rows of the
table use $p\le17$ and the factorial weights directly.

\subsection{Explicit lower bounds from real and from separated conjugates}\label{smd:sec:explicit}
{\raggedright\sources{grh_cm_prime_applications}}

Two elementary constructions give explicit lower bounds for individual discriminants.
The first uses real conjugates and Descartes' rule of signs, the second
\cref{smd:lem:separated}.

\begin{proposition}[Support from coprime factorizations; \stU]\label[proposition]{smd:prop:descartes}
Let $\omega(n)$ be the number of distinct prime divisors of $n$. Then
\[
\begin{aligned}
 s(-4n)&\ge1+2^{\omega(n)-1}&&(n>1),\\
 s(-N)&\ge1+\ceil*{2^{\omega(N)-2}}&&(N>3,\ N\equiv3\pmod4).
\end{aligned}
\]
\end{proposition}

\begin{proof}
There are $2^{\omega(n)-1}$ unordered coprime factorizations $n=uv$ with $u<v$: assign each
full prime power to one of two factors and identify the two orderings. The forms
$[u,0,v]$ are distinct primitive reduced forms of discriminant $-4n$, with CM points
$i\sqrt{v/u}$. All give positive conjugates, since for $z$ on the positive imaginary axis
the formula $j(z)=E_4(z)^3/(q\prod_{m\ge1}(1-q^m)^{24})$, $q=e^{2\pi iz}$, has positive
numerator and denominator by the Eisenstein series expansion. Descartes' rule gives at
most $k-1$ positive roots for a polynomial with $k$ terms. This proves the even bound.

For $N=uv$ odd, again use coprime pairs $u<v$. They are opposite modulo $4$. If $v>3u$,
take $[u,u,(u+v)/4]$. If $v<3u$, take $[(u+v)/4,(v-u)/2,(u+v)/4]$. Equality forces
$(u,v)=(1,3)$, which is excluded. Each triple is integral, reduced and of discriminant
$-N$. It is primitive, because an odd common prime divisor would divide both $u$ and $v$,
and the middle coefficient is odd. Each pair can be recovered from its triple, and the two
types have disjoint shapes. The associated CM points $z=(b+i\sqrt N)/(2a)$ have either
$\Re z=1/2$ or $\abs z=1$. In the first case $-\bar z=z-1$ and in the second
$-\bar z=-1/z$, so modular invariance and the real Fourier coefficients of $j$ show that
$j(z)$ is real. All these conjugates are nonzero because the unique zero of $j$ on the
modular quotient has discriminant $-3$. Of $2^{\omega(N)-1}$ real conjugates at least half
have the same sign. Apply Descartes' rule to $P(X)$ or $P(-X)$.
\end{proof}

\begin{example}\label[example]{smd:ex:840}
At $\Delta=-120$ the complete list of primitive reduced forms is $[u,0,30/u]$ for
$u=1,2,3,5$, and at $\Delta=-840$ it is $[u,0,210/u]$ for $u=1,2,3,5,6,7,10,14$. The
class numbers are $4$ and $8$, and \cref{smd:prop:descartes} gives $s(-120)=5$ and
$s(-840)=9$, the maximum allowed by \eqref{smd:eq:trivial}.
\end{example}

Individual prime discriminants can have arbitrarily large absolute value while giving only
one coprime factor pair, so the factorization bound alone does not grow with $D$.

\begin{proposition}[A certificate with $7{,}958$ conjugates; \stU]\label[proposition]{smd:prop:mersenne}
For $\Delta=1-2^{61}$, every nonzero rational polynomial vanishing at a singular modulus
of discriminant $\Delta$ has at least $7{,}959$ terms.
\end{proposition}

\begin{proof}
Let $D=2^{61}-1$ and $B=\floor{D^{1/4}}=38{,}967$. An integer sieve counts $7{,}958$
integers $1\le a\le B$ all of whose odd prime factors $p$ have $(\Delta/p)=1$. For each such $a$, Hensel
lifting at odd primes and the Chinese remainder theorem show that $\Delta$ is a square
modulo $4a$; the $2$-part follows from $\Delta\equiv1\pmod{2^{61}}$. Every $a$ counted is
coprime to $\Delta$. Hence there is a primitive reduced form with leading coefficient $a$
by \cite[Proposition~3.1(1)]{Bilu2022Trinomials}. For $a<a'\le B$ the exponents
$\pi\sqrt D/a$ and $\pi\sqrt D/a'$ differ by at least $\pi\sqrt D/B^2\ge\pi$. Their
exponentials are both at least $e^{\pi B}$, so the absolute error $2079$ in
\eqref{smd:eq:cusp} changes their magnitudes by less than one percent. The ratio of the
moduli therefore exceeds $(0.99/1.01)e^\pi>16$. Apply \cref{smd:lem:separated} to all
$7{,}958$ moduli. The computation needs neither a numerical value of $j$ nor the
coefficients of the class polynomial.
\end{proof}

\smdcert{the reduced-form lists of \cref{smd:ex:840} and the count of
\cref{smd:prop:mersenne}, with integer arithmetic,
\artifact{grh_cm_prime_applications/verify_cm_support.py}.}

The separated-size method cannot exceed the exponent $1/4$.

\begin{proposition}[Limit of the separated-size method]\label[proposition]{smd:prop:limit}
For $\Delta<-3$, a sequence of conjugates of discriminant $\Delta$ in which each absolute
value is at least sixteen times the preceding one has length $O(D^{1/4})$.
\end{proposition}

\begin{proof}
For a reduced form $[a,b,c]$ with $a\le D^{1/4}$ the conjugate has absolute value within
$2079$ of $\exp(\pi\sqrt D/a)$ by \eqref{smd:eq:cusp}. For sufficiently large $D$ the
largest and smallest of these values at any one $a$ have ratio less than two. Thus this
part supplies at most $\floor{D^{1/4}}$ entries.

For the remaining forms the positive logarithms of the moduli are at most
$\pi D^{1/4}+O(1)$, and the negative logarithms have absolute value $O(\log D)$. For the
latter bound, a reduced point with $\abs b<a$ differs horizontally from either
order-three elliptic corner by at least $1/(2a)$. When $\abs b=a$ the vertical difference
from the closer corner is
\[
 \frac{\abs{D-3a^2}}{2a(\sqrt D+\sqrt3\,a)}\gg D^{-1},
\]
since equality $D=3a^2$ in this case would give the imprimitive multiple of
$[1,\pm1,1]$ or discriminant $-3$, and $a\le\sqrt{D/3}$. The triple zero of $j$ at the
corners and the positive lower bound away from their neighbourhoods give
$\abs j\gg D^{-3}$; see \cite[Lemma~11.6]{Bilu2022Trinomials}. Packing gaps of length
$\log16$ into the resulting interval of logarithms gives the assertion.
\end{proof}

This bound concerns the method, not $s(\Delta)$ itself. The results of the next two
subsections obtain larger exponents by replacing separated archimedean sizes with many
conjugates sharing one residue at a small inert prime.

\subsection{Class numbers and the inert-prime principle}\label{smd:sec:principle}
{\raggedright\sources{grh_cm_sparse_equations; grh_cm_four_terms_and_prime_bias; grh_sparse_singular_moduli}}

\begin{lemma}[Class numbers of orders]\label[lemma]{smd:lem:classnumber}
\leavevmode
\begin{enumerate}[label=\textup{(\alph*)}]
\item Unconditionally, $h(\Delta)=D^{1/2+o(1)}$ as $D\to\infty$ through order
discriminants. The implied lower bound is ineffective.
\item \textup{(\stG)} For all order discriminants,
\[
 \sqrt D/\log\log(3D)\ll h(\Delta)\ll\sqrt D\,\log\log(3D).
\]
\end{enumerate}
\end{lemma}

\begin{proof}
(a) In \eqref{smd:eq:classnumber} the unit index is at most three, and the product lies
between $\varphi(f)/f$ and $\prod_{p\mid f}(1+1/p)$, both $f^{o(1)}$. Siegel's bound
$h(d_K)\gg_\eta\abs{d_K}^{1/2-\eta}$ \cite{Siegel1935Classenzahl} and the elementary upper
bound $h(d_K)\ll\sqrt{\abs{d_K}}\log(2\abs{d_K})$ give the claim.

(b) Write $d=d_K$ and $\chi_d=(d/\cdot)$. By \eqref{smd:eq:classnumber} and the analytic
class number formula, $h(\Delta)$ is $\sqrt D$ times a bounded positive factor times
$L(1,\chi_d)\prod_{p\mid f}(1-\chi_d(p)/p)$. Put $y=(\log(3D))^6$. By
\cref{src:thm:dgrh} the estimate
$\sum_{p\le t}\chi_d(p)\log p\ll\sqrt t\log^2(3Dt)$ holds, and partial summation gives
\[
 \Bigl|\sum_{p>y}\frac{\chi_d(p)}p\Bigr|
 \ll\frac{\log^2(3Dy)}{\sqrt y\log y}
 +\int_y^\infty\frac{\log^2(3Dt)}{t^{3/2}\log t}\,\dd t\ll1 ;
\]
for the last inequality split $\log^2(3Dt)\le2\log^2(3D)+2\log^2t$ and use the choice of
$y$. In the logarithm of the Euler product the terms of prime-power exponent at least two
have bounded absolute sum, and $\sum_{p\mid f,\,p>y}1/p\le\log(2f)/(y\log y)$. Hence
\[
 \log\Bigl(L(1,\chi_d)\prod_{p\mid f}\bigl(1-\chi_d(p)/p\bigr)\Bigr)
 =\sum_{\substack{p\le y\\p\nmid f}}\frac{\chi_d(p)}p+O(1)
 \ge-\log\log y-O(1)
\]
by Mertens' estimate for the sum of reciprocals of primes. The Euler expansion at $1$
follows by taking the limit from the right, which the displayed tail estimate justifies.
This gives the lower bound, and the same argument with the signs reversed gives the upper
bound.
\end{proof}

\begin{lemma}[Inert-prime principle]\label[lemma]{smd:lem:principle}
Let $p\ge5$ be a prime which is inert in $K$ and does not divide $f$, and suppose that at
a place above $p$ at least $h/2$ of the conjugates of discriminant $\Delta$ are units.
Then $s(\Delta)\gg h/p$, with an absolute implied constant.
\end{lemma}

\begin{proof}
By \cref{smd:lem:reduction} the unit conjugates reduce to at most $p/12+O(1)$ nonzero
supersingular invariants, and conjugates with the same reduction lie in one multiplicative
disk of relative depth at least $1$. Some such disk contains $\gg h/p$ distinct roots.
A $t$-term polynomial has at most $\frac{p-1}{p-2}(t-1)\le2(t-1)$ roots there by
\eqref{smd:eq:coarse}.
\end{proof}

\subsection{Growth of the minimum support under Dirichlet GRH}\label{smd:sec:growth}
{\raggedright\sources{grh_cm_sparse_equations; grh_sparse_singular_moduli; grh_cm_four_terms_and_prime_bias; grh_cm_prime_applications}}

\begin{theorem}[Growth of the minimum support; \stG]\label{smd:thm:growth}
Uniformly over imaginary quadratic order discriminants $\Delta=-D$,
\[
 s(\Delta)\gg\frac{h(\Delta)}{(\log(2D))^2}
 \gg\frac{\sqrt D}{(\log(2D))^2\log\log(3D)},
 \qquad
 \lim_{D\to\infty}\frac{\log s(\Delta)}{\log D}=\frac12 .
\]
\end{theorem}

\begin{proof}
First suppose $K\ne\Q(\sqrt{-3})$. On the units modulo $q=6D$ the quadratic characters
attached to $K$ and to $\Q(\sqrt{-3})$ are distinct and nontrivial, since they are induced
by different primitive characters. The simultaneous conditions $\chi_K(p)=-1$,
$\chi_{-3}(p)=1$ specify a coset of a subgroup of index four. With \cref{src:thm:dgrh} as
input, Theorem~1.4 of Lamzouri, Li and Soundararajan \cite{LamzouriLiSoundararajan2015Conditional} yields
a prime $p\ll(\log(2D))^2$ in this coset; small moduli are absorbed in the absolute
constant. This prime is coprime to $6D$, it is inert in $K$, and $p\equiv1\pmod3$. By
\cref{smd:lem:reduction}(iii) the invariant $0$ is ordinary in characteristic $p$, whereas
every conjugate reduces to a supersingular invariant. Hence every conjugate is a unit, and
\cref{smd:lem:principle} gives $s(\Delta)\gg h/p\gg h/(\log(2D))^2$.

For $K=\Q(\sqrt{-3})$ the bound follows from the stronger unconditional estimate of
\cref{smd:prop:eisenstein} below. The second inequality is
\cref{smd:lem:classnumber}(b). The limit follows from the lower bound, from
$s\le h+1$ and from \cref{smd:lem:classnumber}.
\end{proof}

\begin{proposition}[The field $\Q(\sqrt{-3})$; \stU]\label[proposition]{smd:prop:eisenstein}
For $f\ge2$,
\[
 s(-3f^2)\gg\frac{h(-3f^2)}{\log(2f)} .
\]
\end{proposition}

\begin{proof}
Let $\alpha$ be a singular modulus of discriminant $-D=-3f^2$. The complex tori of the
principal lattices $\OO_f\subset\OO_K$ are linked by a cyclic isogeny of degree $f$, and
the second has $j=0$. Pazuki's Theorem~1.1 \cite{Pazuki2016Modular} gives
$h_{\mathrm{Weil}}(\alpha)\le9.204+12\log f$ for the absolute logarithmic Weil height. As
$H_{-D}(0)\ne0$ is an integer,
\[
 \log\abs{H_{-D}(0)}\le h\,h_{\mathrm{Weil}}(\alpha).
\]
At an inert prime $p\nmid D$ the ring class field is unramified, so each nonunit root
contributes at least one to $v_p(H_{-D}(0))$. The total of $\log p$ over the inert primes
at which more than half the roots are nonunits is therefore at most
$2h_{\mathrm{Weil}}(\alpha)$, and the primes dividing $6f$ contribute at most
$\log(6f)$. On the other hand the primes $p\equiv2\pmod3$ up to $X$ have total
$\log p$ asymptotic to $X/2$, by the prime number theorem for this fixed progression.
Taking $X=C\log(2f)$ for a sufficiently large absolute constant $C$ leaves a prime
$p\le X$, inert in $K$ and prime to $6f$, at which at least half the roots are units.
Apply \cref{smd:lem:principle}.
\end{proof}

\subsection{Unconditional growth and near-maximal support}\label{smd:sec:uncond}
{\raggedright\sources{grh_cm_sparse_equations; grh_sparse_singular_moduli; grh_cm_four_terms_and_prime_bias; grh_cm_prime_applications}}

\begin{theorem}[Unconditional growth; \stU]\label{smd:thm:uncond}
For every $\eps>0$ and every order discriminant $\Delta=-D$,
\[
 s(\Delta)\gg_\eps D^{\,1/2-1/(4\sqrt e)-\eps}.
\]
The exponent $1/2-1/(4\sqrt e)$ is approximately $0.348367$. The implied constant is
ineffective. In particular each fixed support size occurs at only finitely many order
discriminants, of arbitrary conductor.
\end{theorem}

\begin{proof}
Enge's Theorem~1.2 gives $\log\abs{H_{-D}(0)}\ll h+\sqrt D\,\log^2(2D)$
\cite{Enge2006Complexity}. Discard $D=3$. At an inert prime $p\nmid D$ at which more than
$h/2$ roots are nonunits, the valuation $v_p(H_{-D}(0))$ is at least $h/2$, so each such
prime costs at least $\tfrac12h\log2$ of this norm. By \cref{smd:lem:classnumber}(a) the
number of these primes is $D^{o(1)}$.

Apply Pollack's Theorem~1.1 \cite{Pollack2015Bounds} to the Kronecker character of
discriminant $-D$, including its imprimitive Euler factors. For each fixed $\delta>0$
there is $c(\delta)>0$ such that, for sufficiently large $D$, more than $D^{c(\delta)}$
primes $p\le D^{1/(4\sqrt e)+\delta}$ have character value $-1$. Every such prime is inert
in $K$ and coprime to $D$. After dropping $2$, $3$ and the $D^{o(1)}$ primes just counted,
at least one remains. \Cref{smd:lem:principle} gives
$s(\Delta)\gg h/D^{1/(4\sqrt e)+\delta}$. Choose $\delta<\eps$ and absorb the loss in
$h=D^{1/2+o(1)}$ into the difference. Siegel's lower bound makes the constant ineffective.
\end{proof}

At two fixed primes the loss disappears almost entirely on an explicit set of
discriminants of positive density.

\begin{theorem}[Near-maximal support in two congruence families; \stU]\label{smd:thm:nearfull}
Let $p\in\set{7,13}$ and let $\Delta$ be an order discriminant with $(\Delta/p)=-1$.
Then $t=s(\Delta)$ satisfies
\[
 h(\Delta)\le t-1+\frac{p-1}{p-2}\bigl(1+\log_pt\bigr),
\]
so that $s(\Delta)=h(\Delta)+1-O(\log(2h(\Delta)))$. The set of order discriminants
satisfying one of the two conditions has relative density $9/13$ among all order
discriminants, counted once each in order of absolute value.
\end{theorem}

\begin{proof}
The condition says that $p$ is inert in $K$ and prime to $f$. By \cref{smd:lem:reduction}
all $h$ conjugates lie in one multiplicative disk of relative depth $1$ about a unit, and
\eqref{smd:eq:digits} with $N=h$ is the displayed inequality. As $t\le h+1$, the additive
loss is $O(\log(2h))$. Order discriminants are the negative integers congruent to $0$ or
$1$ modulo $4$. The nonresidue conditions at $7$ and $13$ hold in proportions $3/7$ and
$6/13$, independently by the Chinese remainder theorem, so their union has relative density
$1-(4/7)(7/13)=9/13$.
\end{proof}

This is a near-equality for the minimum over all multiples of $H_\Delta$, of arbitrary
degree. The density is that of the sufficient congruence conditions.

\subsection{Orders with full support: a criterion and two dyadic families}\label{smd:sec:full}
{\raggedright\sources{grh_cm_sparse_equations; grh_sparse_singular_moduli; grh_transfers_singular_moduli_and_symmetric_powers}}

The equality $s(\Delta)=h(\Delta)+1$ means that \emph{every} choice of $h$ distinct
nonnegative powers of a singular modulus $\alpha$ of discriminant $\Delta$ is a basis of
$\Q(\alpha)$ over $\Q$; equivalently every coefficient of $H_\Delta$ is nonzero and no
multiple of $H_\Delta$, of any degree, has fewer terms than $H_\Delta$. A single disk
containing all conjugates gives a sufficient condition in terms of factorials.

\begin{lemma}[Full-support criterion]\label[lemma]{smd:lem:criterion}
Let $p$ be a prime and $b>1/(p-1)$ a rational number, and suppose that all $h$ conjugates
of discriminant $\Delta$ lie in one multiplicative disk of relative depth $b$ at a place
above $p$. Write $w(n)=bn-v_p(n!)$. If
\begin{equation}\label{smd:eq:G1}
 \max_{0\le u<h}w(u)<\min_{v\ge h}w(v),
\end{equation}
then $s(\Delta)=h+1$. Condition \eqref{smd:eq:G1} is equivalent to
\begin{equation}\label{smd:eq:G2}
 \sum_{t\in I}v_p(t)<b\,\abs I\qquad
 \text{for every interval }I\subset\Z_{>0}\text{ of integers containing }h .
\end{equation}
It holds in particular when the sum of the valuations $v_p(t)$ over each of the intervals
$[h-r+1,h]$ $(1\le r\le h)$ and $[h+1,h+r]$ $(r\ge1)$ is less than $br$.
\end{lemma}

\begin{proof}
Condition \eqref{smd:eq:G1} says that $B_p(h,b)<h$, so by \cref{smd:lem:disk} a polynomial
with at most $h$ terms cannot vanish at $h$ points of the disk; the opposite inequality is
\eqref{smd:eq:trivial}. For $u<h\le v$,
$w(v)-w(u)=\sum_{t=u+1}^{v}(b-v_p(t))$, and the intervals $[u+1,v]$ are exactly those
containing $h$. Under the last condition one has $w(u)<w(h)$ for $u<h$ and
$w(h)<w(v)$ for $v>h$.
\end{proof}

\begin{theorem}[Two dyadic families; \stU]\label{smd:thm:dyadic}
\leavevmode
\begin{enumerate}[label=\textup{(\alph*)}]
\item For every integer $m\ge1$, the discriminant $\Delta=-2^{\,6\cdot7^{m-1}}$ has
$h(\Delta)=2^{\,3\cdot7^{m-1}-2}$ and $s(\Delta)=h(\Delta)+1$.
\item For every integer $k\ge2$, the discriminant $\Delta=-2^{\,2+6\cdot7^{k-1}}$ has
$h(\Delta)=2^{\,3\cdot7^{k-1}-1}$ and $s(\Delta)=h(\Delta)+1$.
\end{enumerate}
For instance $s(-2^{42})=524{,}289$, $s(-2^{44})=1{,}048{,}577$ and
$s(-2^{294})=2^{145}+1$.
\end{theorem}

\begin{proof}
In both cases $\Delta=-4f^2$ is the discriminant of the order of conductor $f=2^n$ in
$\Q(i)$, with $n=3\cdot7^{m-1}-1$ in (a) and $n=3\cdot7^{k-1}$ in (b). By
\eqref{smd:eq:classnumber} its class number is $h=f/2$: the Gaussian field has class
number one, the unit index is two, and $2$ ramifies. The prime $7$ is inert in $\Q(i)$
and coprime to $f$. By \cref{smd:lem:reduction} all $h$ distinct conjugates lie in the
disk $6(1+7\OO)$ of a common unramified completion. Put
$W(n)=n-v_7(n!)=(5n+S_7(n))/6$. By \cref{smd:lem:criterion} with $p=7$ and $b=1$ it
suffices to show that, for every $r\ge1$, the sum of the $7$-adic valuations over each of
the intervals $[h-r+1,h]$ (when $r\le h$) and $[h+1,h+r]$ is strictly less than $r$.

(a) Write $M=7^m$. Because $2^3=1+7$, raising to the power $7^{m-1}$ gives
$4h\equiv1\pmod M$. First take $1\le r<M/5$. The elements of either interval have the
form $h+j$ with $\abs j\le r$. The nonzero integer $1+4j$ has $\abs{1+4j}\le4r+1<M$, so
$v_7(h+j)=v_7(1+4j)$. In either arithmetic progression of $r$ consecutive values of $j$
there are at most $r/7^a+1$ multiples of $7^a$ among the values $1+4j$. Summing over
$a\le\log_7(4r+1)$ bounds the valuation sum by
\[
 r/6+\log_7(4r+1)<r ,
\]
since $4r+1\le5^r$ and $5^6<7^5$. If $r\ge M/5$, direct counting of multiples in the
original interval gives the bound $r/6+\log_7(h+r)$. When $r\le h$,
\[
 \log_7(h+r)\le\log_7(2h)<M/6\le5r/6 ,
\]
where $2^{18}<7^7$ implies $(3/7)\log_72<1/6$. When $r>h$, use $h+r<2r\le5^r$ and
$\log_75<5/6$.

(b) Write $Q=7^k\ge49$. As before $2h=2^n\equiv1\pmod Q$. First let $2r+1<Q$. Then
$v_7(h+j)=v_7(1+2j)$ for $1\le j\le r$, because twice the first number agrees with the
second modulo $Q$ and $0<1+2j<Q$; similarly $v_7(h-j)=v_7(1-2j)$ for $0\le j<r$. Each of
the two sums of valuations is at most $v_7((2r+1)!)$, the only repeated absolute value in
the second sum being $1$. This bound is at most $(2r+1)/6<r$. For $r\ge(Q-1)/2$ use
$S_7(i)\le6(1+\log_7h)$ for $0\le i\le h$: both $W(h+r)-W(h)$ and, for $r\le h$,
$W(h)-W(h-r)$ are at least $(5r-6(1+\log_7h))/6$. Since $\log_7h<n/2=3Q/14$ and
$5(Q-1)/2>6+9Q/7$ for $Q\ge49$, this is positive. The two ranges cover every $r\ge1$.
\end{proof}

The case $k=2$ of (b) has $h=2^{20}$: every coefficient of this class polynomial of degree
$1{,}048{,}576$ is nonzero, and multiplying it by a polynomial of any degree never reduces
its number of terms.
\smdcert{a numerical check of the factorial inequalities for the first four members of
family (a), \artifact{grh_sparse_singular_moduli/checks/check_exact_family.py}.}

\subsection{Full support for Gaussian prime conductors}\label{smd:sec:gaussian}
{\raggedright\sources{grh_transfers_singular_moduli_and_symmetric_powers; grh_sparse_singular_moduli}}

For a rational prime $\ell$ let $j_\ell$ be a singular modulus with endomorphism order
$\Z+\ell\Z[i]$, of discriminant $-4\ell^2$. Its degree is $h=(\ell-\sigma)/2$ for
$\ell$ odd, where $\sigma=\chi_{-4}(\ell)=\pm1$, since $\Z[i]$ has class number one and the
unit index is two.

\begin{theorem}[Gaussian prime conductors; \stU]\label{smd:thm:gaussian}
The primes $\ell$ with $s(-4\ell^2)=h(-4\ell^2)+1$ form a set of lower relative
natural density greater than $0.999818260107372245$ among all primes.
\end{theorem}

The proof places all conjugates in common multiplicative disks at $2$, $3$ and $7$, of
relative depths $4$, $3$ and $2$, and applies \cref{smd:lem:criterion} at each of the
three primes. The exact density of the full-support set is not determined.

\begin{lemma}[Common disks]\label[lemma]{smd:lem:gaussdisks}
Let $\ell\notin\set{2,3,7}$ be prime. In a suitable finite extension of $\Q_p$ with
valuation ring $\OO$, the $h$ conjugates of $j_\ell$ satisfy
\[
 \frac{j}{1728}\in1+7^2\OO\quad(p=7),\qquad
 \frac{j}{1728}\in1+3^3\OO\quad(p=3),\qquad
 \frac{j}{j'}\in1+2^4\OO\quad(p=2)
\]
for any two conjugates $j,j'$ in the last case.
\end{lemma}

\begin{proof}
Every conjugate is the invariant of a quotient of $E_0\colon y^2=x^3-x$ by a cyclic
subgroup of order $\ell$. Indeed its CM lattice corresponds to an invertible ideal $I$ of
the order. Passing from $I$ to $I\Z[i]$ gives index $\ell$, as follows locally from
invertibility. The extended lattice is principal, so the resulting curve has invariant
$1728$, and the dual isogeny gives the desired quotient; see \cite{Sutherland2012Isogeny}
for this lattice description.

\emph{The primes $3$ and $7$.} At $p=3,7$ the curve $E_0$ has good reduction. A finite
unramified extension contains all its $\ell$-torsion, since the torsion of order prime to
$p$ is finite \'etale on the good model. Over its integer ring $\OO$, the quotient by each
cyclic subgroup is an elliptic curve with good reduction, and its reduction is
supersingular because it is isogenous to the supersingular reduction of $E_0$. At $p=7$
use a good integral short equation $y^2=x^3+Ax+B$ of the quotient. The Hasse invariant is
the coefficient of $x^6$ in $(x^3+Ax+B)^3$, which is $3B$. Thus $7\mid B$. The denominator
$4A^3+27B^2$ is a unit, and
\[
 \frac{j}{1728}=\frac{4A^3}{4A^3+27B^2}=1-\frac{27B^2}{4A^3+27B^2}\in1+7^2\OO .
\]
At $p=3$ complete the square in a good integral equation to get
$y^2=x^3+a_2x^2+a_4x+a_6$. The Hasse invariant is the coefficient of $x^2$ in this cubic,
so $3\mid a_2$. Translating the $x$-coordinate by $-a_2/3$ gives an integral short equation
with unit discriminant, and the same formula now gives $j/1728\in1+3^3\OO$, the division
by $1728$ being taken in the local field at $3$.

\emph{The prime $2$.} Let $v=v_2$. For any two cyclic subgroups $H,J$ of $E_0$ of odd
order $\ell$ we claim
\begin{equation}\label{smd:eq:G0}
 v\bigl(j(E_0/H)/j(E_0/J)-1\bigr)\ge4 .
\end{equation}
The integral invariant $1728$ gives $E_0$ potentially good reduction. Take a finite
extension of $\Q_2$ over which $E_0$ has good reduction and which contains
$E_0[3\ell]$. Write $x=r+u^2X$ for the change from the original short equation to a good
integral Weierstrass equation with coordinate $X$. Its discriminant is a unit, while the
original discriminant is $64$, so $v(u^2)=1$. Every nonzero odd-torsion point under
consideration has integral $X$: reduction is injective on prime-to-$2$ torsion, so its
reduction is an affine point of the good model. In particular, for a fixed nonzero
$3$-torsion point $T$ and every nonzero $Q\in H\cup J$ we have $v(x_Q-x_T)\ge1$.

The $3$-division polynomial gives $3x_T^4-6x_T^2-1=0$, whose roots are units. For
$y=5x_T^2-3$ this yields $3y^2-12y-88=0$. The Newton polygon has the single slope $-3/2$,
so $v(y)=3/2$. Since $x_Q=x_T+2z$ with $z$ integral, $v(x_Q^2-x_T^2)\ge2$. For
$g(x)=x(5x^2-3)$ it follows that
\[
 v(g(x_T))=3/2,\qquad v\bigl(g(x_Q)-g(x_T)\bigr)\ge2 .
\]
V\'elu's quotient model is $Y^2=X^3+A_HX+B_H$, where
\[
 A_H=-1-10\!\!\sum_{Q\in(H\setminus\{0\})/\{\pm1\}}\!\!(3x_Q^2-1),\qquad
 B_H=-14\!\!\sum_{Q\in(H\setminus\{0\})/\{\pm1\}}\!\!g(x_Q);
\]
these are the short-equation specialization of the formulas of
\cite[Theorem~1]{Miret2007Generalization}. In particular $A_H$ is a unit congruent to $-1$
modulo $2$ and $v(B_H)\ge5/2$. The sums for $H$ and $J$ have the same number
$(\ell-1)/2$ of terms, all congruent to $g(x_T)$ modulo $4$ in the second sum. Thus
$v(B_H-B_J)\ge3$ and $v(B_H+B_J)=v(B_H-B_J+2B_J)\ge3$, which gives
$v(B_H^2-B_J^2)\ge6$.

Put $s_H=27B_H^2/(4A_H^3)$. Then $v(s_H)\ge3$ and $v(s_H-s_J)\ge4$: replacing $B_H^2$ by
$B_J^2$ costs at least four in valuation after division by $4$, while
$v(A_H^{-3}-A_J^{-3})\ge1$. The identity $j(E_0/H)=1728/(1+s_H)$ proves
\eqref{smd:eq:G0}.
\end{proof}

Fixing any one of the quotient invariants $j_*$ therefore puts all $h$ conjugates in
$j_*(1+16\OO)$; the centre is algebraic, which \cref{smd:lem:disk} allows.

\begin{remark}[Sharpness of the depth at $2$]\label[remark]{smd:rem:depth2}
The depth four in \eqref{smd:eq:G0} is sharp when the assertion ranges over all odd prime
conductors. For $\ell=3$ let $Y=j/1728-1$. Substituting $A=9-30x_T^2$,
$B=-14x_T(5x_T^2-3)$ into the formula for $j$ and eliminating by
$3x_T^4-6x_T^2-1=0$ gives the quadratic $27Y^2-2399040Y-18593344$. Its discriminant has
$2$-adic valuation eight, so its two roots differ by an element of valuation four, and
both $1+Y$ are units. The corresponding pair attains equality in \eqref{smd:eq:G0}.
\smdcert{the elimination, and the corresponding check for the isogenies of degree five,
\artifact{grh_transfers_singular_moduli_and_symmetric_powers/quotient_check.py}.}
\end{remark}

\begin{proof}[Proof of \cref{smd:thm:gaussian}]
Discard $\ell=2,3,7$. For $(p,b)\in\set{(2,4),(3,3),(7,2)}$ say that $\ell$ \emph{passes
the test at $p$} if \eqref{smd:eq:G2} holds for $h=(\ell-\sigma)/2$ with these $p$ and
$b$. By \cref{smd:lem:gaussdisks,smd:lem:criterion}, a prime passing at least one of the
three tests has $s(-4\ell^2)=h+1$.

\emph{Finite densities.} Fix $\sigma$. For lengths $\abs I\le R$, and omitting the
finitely many $h<R$, the test \eqref{smd:eq:G2} depends only on $h\bmod p^{bR}$. For an
odd $p$ the admissible classes of $h$ modulo $p^a$ are those with $p\nmid2h+\sigma$. For
$p=2$ they are exactly the even classes: $\ell=2h+\sigma\equiv\sigma\pmod4$, and all such
classes correspond to invertible classes of $\ell$ modulo $2^{a+1}$ with that sign. In all
three cases each admissible class has conditional density $1/\varphi(p^a)$, by the prime
number theorem in arithmetic progressions. Hence the set of primes of sign $\sigma$
satisfying \eqref{smd:eq:G2} for all intervals of length at most $R$ has a relative
density $d_{p,R}$. Reflection $h\mapsto-h$ and reversal of intervals show that $d_{p,R}$
is the same for either $\sigma$.

\emph{The tail.} If an interval of length $r$ violates \eqref{smd:eq:G2} and its greatest
valuation is $\mu$, counting multiples of $p^a$ for $1\le a\le\mu$ bounds the sum of
valuations by $r/(p-1)+\mu$. Consequently some $h+j$ with $\abs j\le r-1$ is divisible by
\[
 q_r=p^{\ceil{(b-1/(p-1))r}} .
\]
For fixed $r$ the conditional upper density of this event is at most
$u_r=(2r-1)/\varphi(q_r)$. To sum these bounds, let $X$ be the cutoff for $\ell$. For
$q_r\le\sqrt X$ the Brun--Titchmarsh inequality \cite{Montgomery1973Large} bounds the
number of primes in each relevant invertible residue class by
$O(X/(\varphi(q_r)\log X))$, with an absolute constant. The modulus is $q_r$ for odd $p$
after dropping the condition modulo $4$, and $2q_r$ for $p=2$, because
$\ell=2h+\sigma$; in the latter case $2q_r\le2\sqrt X$ still gives
$\log(X/(2q_r))\gg\log X$. Since only $\ell=p$ can occupy a noninvertible class, dropping
that prime and dividing by the number of primes of each sign gives a uniform tail bounded
by an absolute multiple of $\sum_{r>T}u_r$. For $q_r>\sqrt X$, a positive multiple of
$p^{\floor{\log_p\sqrt X}}$ occurs among the $h+j$. The inequality
$(b-1/(p-1))r\le\log_p(h+r)$ forces $r=O(\log X)$: when $r>h$ it fails, since
$h+r<2r$ and $2r\le2^r<p^{(b-1/(p-1))r}$, and otherwise $h+r\ll X$. There are
$O(\sqrt X\log X)$ possibilities for a multiple and an offset, which is negligible
compared with the number of primes. Taking limits at fixed $T$ and then $T\to\infty$
proves that the density $d_p$ of primes passing the test at $p$ exists and that
\begin{equation}\label{smd:eq:G3}
 0\le d_{p,R}-d_p\le\sum_{r>R}u_r .
\end{equation}

\emph{Exact certificate.} The finite density is computed as follows. Start with the
$p-1$ admissible classes modulo $p$. In a class modulo $p^a$, evaluate all relevant
valuations truncated at $a$. If any interval sum already reaches $br$, discard the class.
If no valuation is truncated, accept the class. Otherwise visit all $p$ lifts. The process
terminates by depth $bR$, since a truncated term at that depth already violates a length
at most $R$ containing that term and the centre. To bound \eqref{smd:eq:G3}, group its
terms in successive blocks of length $p-1$: for $r\ge R+1$ the ratio
\[
 \frac{u_{r+p-1}}{u_r}=\frac{2r+2p-3}{2r-1}\,p^{-(b(p-1)-1)}
\]
is decreasing, and a geometric sum bounds the whole tail by its first block. All
operations use rational arithmetic.
\begin{center}
\begin{tabular}{cclc}
\toprule
$p$ & $R$ & $d_{p,R}$ & Tail upper bound\\
\midrule
2&25&$65963232382760376856543/75557863725914323419136$ &$3.879\cdot10^{-22}$\\
3&16&$68736524642595367189/72945992754341572806$ &$1.694\cdot10^{-19}$\\
7&10&$9528160153068835/9770481587462694$ &$5.089\cdot10^{-17}$\\
\bottomrule
\end{tabular}
\end{center}

\emph{The union.} At fixed sign the three finite congruence tests are independent, by the
Chinese remainder theorem and the prime number theorem in progressions. Passing to the
limits using \eqref{smd:eq:G3}, the density of primes passing at least one test is
$d=1-(1-d_2)(1-d_3)(1-d_7)$. It is the same for both signs and satisfies
\[
 0.9998182601073722452943<d<0.9998182601073722456678 .
\]
Every prime counted by $d$ gives full support.
\end{proof}

\smdcert{\artifact{grh_transfers_singular_moduli_and_symmetric_powers/check_density.py},
with output
\artifact{grh_transfers_singular_moduli_and_symmetric_powers/checks/check_density.txt}.}

\begin{example}[Individual conductors]\label[example]{smd:ex:conductors}
For $\ell=197$ the degree is $98$. The one-term interval $\set{98}$ fails
\eqref{smd:eq:G2} at $7$, but the test at $3$ is passed: the singleton has valuation
$0$, and every interval of length $r\ge2$ through $98$ has valuation sum at most
$r/2+\log_3(98+r)<3r$, an inequality which holds at $r=2$ and persists since
$3^{5r/2}$ grows faster than $98+r$ at every step. Thus any $98$ distinct nonnegative
powers of $j_{197}$ form a basis of $\Q(j_{197})$.

For $\ell=107$ the degree is $54$. Here $3^3\mid54$ rules out \eqref{smd:eq:G2} at $3$,
while $7\nmid54$, and for $r\ge2$ one has $r/6+\log_7(54+r)<2r$. Every $54$ distinct
powers of $j_{107}$ form a basis.

The conductor $15877$ shows the contribution of the prime $2$. It is prime and
$h=7938=2\cdot3^4\cdot7^2$, so both odd-prime tests fail on the singleton $\set h$. The
largest sums of $2$-adic valuations of intervals through $7938$ of lengths $1,2,3,4$ are
$1,1,9,9$, each smaller than four times the length. For every $r\ge5$ the sum is at most
$r+\log_2(7938+r)<4r$, since $7938+5<2^{15}$ and the inequality persists on increasing
$r$. Thus any $7938$ distinct nonnegative powers of this singular modulus form a basis.
\end{example}

\begin{remark}[Conductors up to $100{,}000$]\label[remark]{smd:rem:missed}
Among the $9{,}589$ primes $\ell\le100{,}000$ other than $2,3,7$, exactly $9{,}587$ pass
at least one test, and the two odd primes alone certify $9{,}575$; for a given $h$ the
direct check of \eqref{smd:eq:G1} can stop as soon as $n(b-1/(p-1))>bh$. The first conductor missed by all three tests is $\ell=42{,}337$,
with $h=21{,}168=16\cdot27\cdot49$: the singleton $\set h$ defeats all three tests. The
other is $\ell=84{,}673$, with $h=42{,}336$. The exact least support in these two cases is
not determined.
\smdcert{\artifact{grh_transfers_singular_moduli_and_symmetric_powers/checks/conductors.txt}.}
\end{remark}

\subsection{The discriminant \texorpdfstring{$-87$}{-87}: powers and residue patterns}\label{smd:sec:powered}
{\raggedright\sources{grh_transfer_bridges_2026-11-25; padic-six-term; singular-modulus-support}}

For an algebraic number $\gamma$ write $s(\gamma)$ for the least number of nonzero terms
of a nonzero rational polynomial vanishing at $\gamma$. At discriminant $-87$ one knows
$6\le s(-87)\le7$ (\cref{smd:prop:sextics}), and no five of its conjugates lie in a disk deep enough for a
disk certificate (\cref{smd:rem:disklimit}). The higher-depth case of
\cref{smd:thm:confluent} settles the relations in one residue class of exponents, and more
generally every relation whose pattern of residues passes a finite rank test.

\begin{theorem}[A powered singular modulus; \stU]\label{smd:thm:powered}
Let $\alpha$ be any singular modulus of discriminant $-87$. Then $s(\alpha^{63t})=7$ for
every integer $t\ge1$. The exponents and rational coefficients of the annihilating
polynomials are unrestricted.
\end{theorem}

\begin{proof}
The ring class polynomial $H_{-87}$ has degree six; it is the certified polynomial of
\cref{smd:sec:sextics}. In $A=\Z_{(2)}[X]/(H_{-87})$, with $x$ the class of $X$, exact
reduction gives
\[
 x^{63}\equiv1+16(x+x^2+x^3)\pmod{32A}.
\]
Thus $x^{63}=1+16D_0$ with $D_0\bmod2=x+x^2+x^3$. The six columns $1,D_0,\dots,D_0^5$,
represented with entries $0,1$ in the power basis, have integer determinant $-1$. They are
independent modulo $2$, and $6<2^4$. By \cref{smd:thm:confluent} with $p=2$, $\sigma=4$ and
$m=63$, any six distinct powers of $x$ with exponents divisible by $63$ are linearly
independent. A rational relation for $\alpha^{63t}$ with at most six terms would select at
most six such powers, and a smaller dependent set could be enlarged to six distinct
exponents, a contradiction. On the other hand $\alpha^{63t}$ lies in a six-dimensional
$\Q$-algebra, so seven of its powers are dependent. The modular power was verified both by
polynomial multiplication and by an independently written basis-shift recurrence, and the
determinant with integer arithmetic.
\end{proof}

The companion statement $s(\beta^{182t})=7$ for a singular modulus $\beta$ of discriminant
$-152$, originally proved by the same method with the congruence
$x^{182}\equiv26+9(2x^2+x^3+2x^5)\pmod{27}$ in $\Z_{(3)}[X]/(H_{-152})$, now follows from
\cref{smd:thm:seven}: a relation for $\beta^{182t}$ with at most six terms is a relation for
$\beta$ with at most six terms, and $\beta^{182t}$ lies in a six-dimensional algebra.

\begin{proposition}[Residue-pattern certificate at $-87$; \stU]\label[proposition]{smd:prop:eightyseven}
Let $\F_2[X]/(X^6+X^5+X^2+X+1)$ be the reduction of $\Z_{(2)}[X]/(H_{-87})$, in which the
class of $X$ has order $63$, and put $\ell=X+X^2+X^3$. Let $e_1<\dots<e_6$ be integers, write
$e_i-e_1=a_i+63q_i$ with $0\le a_i<63$, and let $\mu_a$ be the number of $i$ with $a_i=a$.
If the six columns
\[
 X^a\ell^{\,j},\qquad 0\le a<63,\quad 0\le j<\mu_a,
\]
are linearly independent over $\F_2$, then the powers $\alpha^{e_1},\dots,\alpha^{e_6}$ of a
singular modulus $\alpha$ of discriminant $-87$ are linearly independent over $\Q$. The test
depends only on the multiset of residues $a_i$; among the $\binom{67}5=9{,}657{,}648$
multisets of six residues modulo $63$ containing $0$, exactly $3{,}892{,}897$ pass it.
\end{proposition}

\begin{proof}
Multiplication by $\alpha^{-e_1}$ reduces to nonnegative exponents with $e_1=0$. In the
notation of the proof of \cref{smd:thm:powered}, $x^{63}=1+16D_0$ with $D_0\equiv\ell$
modulo $2$, each residue block has at most six elements and $6<2^4$, so
\cref{smd:thm:confluent} with $p=2$, $\sigma=4$, $m=63$ and $a=1$ applies; its columns
$x^aD_0^{\,j}$ reduce to $X^a\ell^{\,j}$. A full-rank residue matrix therefore certifies all
lifts of that residue multiset at once. For a single residue the six powers of $\ell$ have
bit masks $1,14,51,7,42,48$ and are independent, which recovers \cref{smd:thm:powered}. The
count is an exhaustive enumeration of the weakly increasing residue multisets with smallest
entry $0$, to which every pattern can be translated by a common exponent shift.
\end{proof}

A failed rank test is inconclusive, and the remaining $5{,}764{,}751$ residue multisets are
not covered by this certificate.
\smdcert{\artifact{grh_transfer_bridges_2026-11-25/checks/confluent_certificates.json}, where the
polynomials coincide with those of
\artifact{grh_sparse_singular_moduli/checks/five_small_witnesses.json}; checker
\artifact{grh_transfer_bridges_2026-11-25/checks/check_confluent_refinement.py}, with
output \artifact{grh_transfer_bridges_2026-11-25/evidence/confluent_replay.txt}; congruence
modulo $32$ and the multiset count \artifact{padic-six-term/verify_87.py}, with output
\artifact{padic-six-term/verification_87.txt}.}

\subsection{Open questions}\label{smd:sec:open}
{\raggedright\sources{grh_sparse_singular_moduli; grh_cm_sparse_equations; grh_transfers_singular_moduli_and_symmetric_powers; grh_transfer_bridges_2026-11-25; singular-modulus-support; singular_moduli_seven_terms_2026-11-26; padic-six-term}}

\begin{question}[Six terms]\label{smd:q:six}
Does every singular modulus of degree $h\ge6$ require at least seven terms, that is, does
$h(\Delta)\ge6$ imply $s(\Delta)\ge7$? In particular, is $s(-87)=7$? For $-152$ and nine
further sextic orders, and for $-4075$, the answer is yes
(\cref{smd:thm:seven,smd:ex:4075}), and four orders of degree seven have $s=8$. At $-87$,
\cref{smd:thm:powered,smd:prop:eightyseven} exclude six-term relations for $3{,}892{,}897$
of the $9{,}657{,}648$ anchored residue patterns modulo $63$, including the relations whose
exponents lie in one residue class. A global statement would need, in addition to a finite
residue cover for the low-degree orders, a large-discriminant reduction of the kind supplied
by \cref{smd:prop:global} for five terms. More generally, is there a uniform larger
threshold?
\end{question}

\begin{question}[A fixed fraction of the class number]\label{smd:q:fraction}
Is $s(\Delta)\ge c\,h(\Delta)$ for an absolute constant $c>0$ and every order? The
argument of \cref{smd:thm:growth} loses a factor of order $p$ when the conjugates occupy
different supersingular residues, and \cref{smd:thm:nearfull} gives a ratio tending to
one only on the family of density $9/13$.
\end{question}

\begin{question}[Gaussian prime conductors]\label{smd:q:density}
Is the relative density of primes $\ell$ with $s(-4\ell^2)=h+1$ equal to one? What is
$s(-4\ell^2)$ for $\ell=42{,}337$ and $\ell=84{,}673$? At these conductors all three
first-difference tests fail at the central integer. A tied coefficient valuation leaves
room for a nonzero generalized Vandermonde determinant, and its leading residue for the
actual CM points is a more precise target than the valuations of the factorials alone.
\end{question}

\begin{question}[Other CM isogeny classes]\label{smd:q:isogeny}
Can the common-disk argument of \cref{smd:lem:gaussdisks} be extended to other CM isogeny
classes? The proof uses a model with good reduction over an unramified field, an isogeny
of degree prime to the residue characteristic, and explicit control of the local
coordinate of every quotient; the extra depth at $3$ comes from normalizing by $1728$
before taking logarithms.
\end{question}
